%%==============================================================================
%%  PAPER 9 -- LOCKED BLUEPRINT
%%  Direction-dependent band gaps in 2D anisotropic strain-gradient elastic media
%%
%%  Target: Q1 journal, Impact Factor >= 5  (primary: Int. J. Mechanical Sciences)
%%
%%  Compile:  pdflatex Paper9_Blueprint.tex  (run twice for TOC + refs)
%%
%%  This file is BOTH the blueprint document AND the manuscript skeleton.
%%  Sections marked  [BLUEPRINT]  are planning annotations -- delete them when
%%  you convert this file into the actual submission.
%%==============================================================================
\documentclass[11pt,a4paper]{article}

\usepackage[margin=1.85cm,top=2.0cm,bottom=2.0cm]{geometry}
\usepackage{booktabs,longtable,tabularx,array,multirow}
\usepackage{xltabular}
\usepackage[table]{xcolor}
\usepackage{enumitem}
\usepackage{titlesec}
\usepackage{fancyhdr}
\usepackage{amsmath,amssymb,bm}
\usepackage{ragged2e}
\usepackage[hidelinks]{hyperref}
\usepackage{colortbl}

%% ---- colours ---------------------------------------------------------------
\definecolor{lockblue}{HTML}{1F3864}
\definecolor{accent}{HTML}{2E74B5}
\definecolor{warn}{HTML}{C00000}
\definecolor{ok}{HTML}{1E7B34}
\definecolor{softgrey}{HTML}{F2F2F2}
\definecolor{softblue}{HTML}{DEEAF6}
\definecolor{softgreen}{HTML}{E2EFDA}
\definecolor{softamber}{HTML}{FFF2CC}
\definecolor{softred}{HTML}{FBE4E4}
\definecolor{addcol}{HTML}{7030A0}
\definecolor{strcol}{HTML}{C55A11}
\definecolor{corcol}{HTML}{0F7B6C}

%% ---- headings --------------------------------------------------------------
\titleformat{\section}{\Large\bfseries\color{lockblue}}{\thesection}{0.7em}{}
\titleformat{\subsection}{\large\bfseries\color{accent}}{\thesubsection}{0.6em}{}
\titleformat{\subsubsection}{\normalsize\bfseries}{\thesubsubsection}{0.6em}{}
\titlespacing*{\section}{0pt}{14pt}{6pt}
\titlespacing*{\subsection}{0pt}{10pt}{4pt}

\setlist[itemize]{leftmargin=1.3em,itemsep=1.5pt,topsep=3pt}
\setlist[enumerate]{leftmargin=1.6em,itemsep=1.5pt,topsep=3pt}

\pagestyle{fancy}\fancyhf{}
\fancyhead[L]{\footnotesize\color{accent}Paper 9 --- Locked Blueprint}
\fancyhead[R]{\footnotesize\color{accent}\thepage}
\renewcommand{\headrulewidth}{0.4pt}

%% ---- helper macros ---------------------------------------------------------
\newcommand{\bp}[1]{\textcolor{accent}{\textsf{\footnotesize [#1]}}}
\newcommand{\ADDED}{{\color{addcol}\textbf{[ADDED]}}}
\newcommand{\STR}{{\color{strcol}\textbf{[STRENGTHENED]}}}
\newcommand{\COR}{{\color{corcol}\textbf{[CORRECTED]}}}
\newcommand{\chk}{$\square$~}
\newcommand{\YES}{\textcolor{ok}{\textbf{\checkmark}}}
\newcommand{\NO}{{\color{warn}\textbf{X}}}
\newcommand{\PART}{\textcolor{warn}{$\sim$}}
\newcolumntype{L}[1]{>{\RaggedRight\arraybackslash}p{#1}}
\newcolumntype{Y}{>{\RaggedRight\arraybackslash}X}

\setlength{\parskip}{3pt}
\setlength{\parindent}{0pt}
\sloppy
\emergencystretch=2.5em
\hfuzz=2pt
\renewcommand{\arraystretch}{1.18}

%%==============================================================================
\begin{document}

%%------------------------------------------------------------------------------
\begin{titlepage}
\centering
\vspace*{0.6cm}
{\color{accent}\rule{\textwidth}{1.2pt}}\\[0.5cm]
{\Huge\bfseries\color{lockblue} PAPER 9}\\[0.35cm]
{\Large\bfseries LOCKED BLUEPRINT}\\[0.9cm]
{\color{accent}\rule{\textwidth}{0.6pt}}\\[1.1cm]

{\large\bfseries Direction-dependent band gaps and iso-frequency contours in\\[0.25cm]
two-dimensional strain-gradient elastic media with anisotropic\\[0.25cm]
ellipsoidal microstructure: a Bloch--Floquet $C^1$ finite element analysis}\\[1.3cm]

\begin{minipage}{0.86\textwidth}
\centering\footnotesize
\begin{tabular}{@{}>{\bfseries}p{3.5cm} p{9.6cm}@{}}
\textbf{Document type} & Locked blueprint + manuscript skeleton \\
\textbf{Version} & 1.5 --- LOCK; amendment \textbf{A2}: external-validation evidence routes (quantitative / graphical / not-validated; defined in Section~13) with the \textbf{quantitative criterion and thresholds unchanged} ($\le 2\%$; $0.5\%$ classical target). v1.5 is derived from v1.4 (sha256 \texttt{2ae0b1e8f37e10a0685a0b0ffd94e695bd62e3b4da6a02052a76190fc0acb638}) \emph{plus} the A2 blocks itemised in \texttt{paper9/audit/BLUEPRINT\_V1\_5\_GRAPHICAL\_VALIDATION\_AMENDMENT.md}; v1.4 is preserved byte-identical and records its own provenance from v1.3 (\texttt{ca71b91aba4ca4abe9f157eb595ac8c2f709d838957a08ea0dce7160685dbf9f}). The \S5.7 Rule R-fit text (pre-declared factor $F = 3$, declared zero $\max(s_i, 10^{-15})$) is retained unchanged. \\
\textbf{Date} & 22 September 2026 \\
\textbf{Principal Investigator} & Dr.\ Vipin Gupta, Dept.\ of Mathematics, Gurugram University \\
\textbf{Funding} & Haryana State Research Fund, Sanction Order No.\ 14/66-2025 SPD/HSHEC/HSRF \\
\textbf{Primary target journal} & International Journal of Mechanical Sciences (Q1, IF 11.4) \\
\textbf{Backup journals} & J.\ Sound and Vibration (5.8) $\cdot$ Composite Structures (7.8) $\cdot$ Appl.\ Math.\ Modelling (5.1) \\
\textbf{Planned length} & 30--34 pages, 71 main-text equations + Appendices A--B, 13 figures, 6 tables \\
\textbf{Planned effort} & 14 weeks \\
\end{tabular}
\end{minipage}

\vfill
{\color{accent}\rule{\textwidth}{0.6pt}}\\[0.35cm]
{\footnotesize\itshape Companion files: \texttt{ANCHOR\_DATA\_SHEET.md} (validation sources) $\cdot$
\texttt{MASTER\_ROADMAP.md} (asset library) $\cdot$
\texttt{anchors/} (downloaded 2023/2024 benchmark PDFs)}\\[0.3cm]
{\color{accent}\rule{\textwidth}{1.2pt}}
\end{titlepage}

%%------------------------------------------------------------------------------
\tableofcontents
\newpage

%%==============================================================================
\section{LOCK SHEET --- read this first}
%%==============================================================================

\subsection{The five locked decisions}

{\small
\begin{tabularx}{\textwidth}{@{}>{\bfseries}p{0.5cm} L{3.6cm} Y@{}}
\toprule
\# & Decision & Locked value \\
\midrule
D1 & Analysis class & \textbf{Bloch--Floquet band structure} of an infinite 2D periodic medium.
       \emph{Not} a static solve, \emph{not} a bounded-structure eigenproblem. \\
D2 & Constitutive model & Mindlin Form-II strain-gradient elasticity with an
       \textbf{anisotropic ellipsoidal characteristic-length tensor} \emph{plus}
       \textbf{micro-inertia}. Linear, isothermal, non-piezoelectric. \\
D3 & Exactly one new ingredient & \textbf{Dynamics + periodicity + micro-inertia.}
       Explicitly \emph{excluded}: fractional/Rabotnov memory, functionally graded
       material, flexoelectricity, piezoelectricity, thermal field, geometric
       nonlinearity, 3D. \\
D4 & Novelty hook & \textbf{Direction-dependent band gaps.} Rotating the ellipsoidal
       microstructure steers which propagation directions are stopped. \STR\ Within the scope of our review (Table~1), and to the best of our knowledge, all published gradient-elastic band-gap work is 1D or isotropic. \\
D5 & Validation & \COR\ \textbf{Five-layer validation and verification framework.} Layers 1--4 are external validation (published benchmarks, closed-form analytics, independent published method); Layer 5 is internal numerical verification (eight-test suite + mesh convergence/resolution floor). The three published anchors (\textbf{2023, 2024, 2026}) sit in Layers 1, 2 and 4. \STR\ The benchmark configurations, exact parameters, reproduced curves and error tables appear \emph{in the submitted manuscript itself}, never only in supplementary or internal files. Hard gate: the paper is \emph{not} submitted until Layers 1, 2a and 2b agree to $\le 2\%$ (validation of the physical formulation); Layer 5 verifies the numerical implementation. \\
\bottomrule
\end{tabularx}
}

\subsection{Contribution statement \bp{goes verbatim into the Introduction}}

\begin{enumerate}[label=\textbf{C\arabic*.}]
\item \textbf{A two-dimensional Bloch--Floquet framework for anisotropic strain-gradient
elasticity.} The characteristic-length tensor is derived from the second moments of a
rotated ellipsoidal nonlocal averaging domain, so that microstructural aspect ratio
$\mathrm{AR}$ and orientation $\theta$ enter the dispersion problem as two independent,
physically interpretable design variables.

\item \textbf{Micro-inertia is included as an independent constitutive parameter
$\ell_{\mathrm{i}}$, and its necessity is demonstrated.} \STR\ Without it the phase velocity
grows without bound as $k\to\infty$: Appendix~A proves this asymptotically and \S7.4 confirms it numerically, so the model without micro-inertia loses bounded wave speed and is physically inadmissible. Causality-type language is used only with both supports, never from a numerical plot alone; the bounded/unbounded transition is quantified rather than asserted.

\item \textbf{A $C^1$ Bogner--Fox--Schmit Bloch finite element in which the complex
phase factor is applied consistently to the value \emph{and} to all first- and
mixed-derivative degrees of freedom.} This is required by the fourth-order operator
and, \STR\ to the best of our knowledge, has not been documented for gradient-elastic band-structure analysis (Table~1).

\item \textbf{Direction-dependent band gaps, quantified.} Complete and partial band-gap
widths are mapped on the $(\theta,\mathrm{AR})$ plane, and iso-frequency contours and
group-velocity fields are used to show that the stop-band direction can be steered by
microstructural orientation alone --- a manufacturing-controllable knob.
\end{enumerate}

\subsection{Uniqueness declaration \bp{answers the ``isn't this an extension?'' risk}}

\begin{tabularx}{\textwidth}{@{}L{3.1cm} Y Y@{}}
\toprule
Attribute & Prior static gradient-elastic FE work & \textbf{This paper} \\
\midrule
Domain & Bounded structure (bar, rectangle, cuboid) & Infinite periodic medium: unit cell + Bloch phase \\
Primary variable & Load / displacement & Wavenumber $\bm{k}$ in the Brillouin zone \\
Observable & Stiffening ratio, displacement, boundary-layer thickness, frequency $f_1$ & Band-gap edges and widths, iso-frequency contours, group velocity \\
Constitutive & No gradient inertia (not needed in statics) & \textbf{Micro-inertia $\ell_{\mathrm i}$ present and essential} \\
Boundary conditions & Real periodic tying, or clamped/traction & \textbf{Complex Bloch phase on value \emph{and} derivative DOFs} \\
Spectrum & Discrete eigenvalues & Continuous bands separated by gaps \\
Physics & Size-dependent stiffening & Wave filtering / Bragg scattering + microstructure-induced dispersion \\
\bottomrule
\end{tabularx}

\medskip
{\color{warn}\textbf{Self-containment rule (mandatory).}}
The five existing manuscripts are \emph{not yet published}. This paper therefore
\textbf{re-derives every constitutive and discretisation result from first principles}
and cites nothing of ours as a source. It must be publishable before, after, or
independently of any other manuscript in the series. No forward dependency exists.

\subsection{Journal target}

{\small
\begin{tabularx}{\textwidth}{@{}c L{3.6cm} c c Y@{}}
\toprule
Rank & Journal & IF & Quartile & Fit \\
\midrule
\cellcolor{softgreen}1 & \cellcolor{softgreen}International Journal of Mechanical Sciences (Elsevier) & \cellcolor{softgreen}11.4 & \cellcolor{softgreen}Q1 & \cellcolor{softgreen}Publishes exactly this class: strain-gradient continua, metamaterial band gaps, $C^1$/higher-order FEM. Highest upside. \\
2 & Journal of Sound and Vibration (Elsevier) & 5.8 & Q1 & Wave propagation, dispersion, vibration isolation. Matches the stated IF~5 bar. \\
3 & Composite Structures (Elsevier) & 7.8 & Q1 & Microstructured/architected media, size effects. \\
4 & Applied Mathematical Modelling (Elsevier) & 5.1 & Q1 & Rewards the formulation + verification rigour. \\
5 & Waves in Random and Complex Media (T\&F) & $\approx$2.7 & Q1/Q2 & \textbf{Strategic value:} both validation anchors (A and C) and two of our own papers are here. Below the IF bar, so kept as fallback only. \\
\bottomrule
\end{tabularx}
}

\medskip
\bp{All four top targets are Elsevier. Use the \texttt{elsarticle} class; the blueprint below is identical for all of them. Journal-specific deltas are listed in Section~\ref{sec:submission}.}

\newpage
%%==============================================================================
\section{FRONT MATTER}
%%==============================================================================

\subsection{Title}

\textbf{Locked (primary):}
\begin{quote}\itshape
Direction-dependent band gaps and iso-frequency contours in two-dimensional
strain-gradient elastic media with anisotropic ellipsoidal microstructure:
a Bloch--Floquet $C^1$ finite element analysis
\end{quote}

\textbf{Alternate A (shorter, for JSV):}
\begin{quote}\itshape
Orientation-tunable elastic band gaps in two-dimensional Mindlin Form-II
strain-gradient media with micro-inertia
\end{quote}

\textbf{Alternate B (method-forward, for IJMS):}
\begin{quote}\itshape
A Bloch--Floquet $C^1$ finite element for anisotropic strain-gradient elasticity:
direction-dependent band gaps and wave steering in microstructured media
\end{quote}

\subsection{Author and funding block}

\begin{itemize}
\item \textbf{Authors / affiliations / emails / ORCID} --- to be finalised. \textcolor{warn}{Do not leave blank} (this was the defect in manuscript FEM\_3).
\item \textbf{Corresponding author} with full postal address, email, ORCID.
\item \textbf{Funding statement (mandatory verbatim):}
\begin{quote}\footnotesize\itshape
``This work was supported by the Haryana State Research Fund, Haryana State Higher
Education Council, Sanction Order No.\ 14/66-2025 SPD/HSHEC/HSRF, Session 2025--2026.''
\end{quote}
\textcolor{warn}{Missing in FEM\_1, FEM\_2, FEM\_3, GK, QC. Must appear here.}
\item \textbf{CRediT author statement} --- required by all four target journals.
\item \textbf{Declaration of competing interest} --- required.
\item \textbf{Data availability} --- ``The MATLAB/Python code and all computed band-structure datasets supporting the conclusions of this article are available at \texttt{[repository URL]} (DOI: \texttt{[Zenodo DOI]}).'' \bp{this is also evidence for project outcome X3}
\end{itemize}

\subsection{Abstract --- sentence-by-sentence plan \bp{target 230--260 words}}

{\small
\begin{tabularx}{\textwidth}{@{}>{\bfseries}p{0.7cm} L{2.7cm} >{\centering\arraybackslash}p{1.0cm} Y@{}}
\toprule
S\# & Function & Words & Content to state \\
\midrule
1 & Context & 20 & Microstructured solids filter elastic waves; at wavelengths comparable to the microstructure, classical elasticity fails. \\
2 & Theory gap & 25 & \STR\ Gradient-elastic band-gap studies to date have, to the best of our knowledge, been one-dimensional or isotropic, so the role of microstructural \emph{orientation} is unknown. \\
3 & What we do & 30 & A 2D Bloch--Floquet $C^1$ FE framework for Mindlin Form-II strain-gradient elasticity with an anisotropic ellipsoidal characteristic-length tensor and micro-inertia. \\
4 & Method novelty & 25 & Complex Bloch phase applied consistently to value and derivative DOFs; Hermitian reduced eigenproblem over the irreducible Brillouin zone. \\
5 & Validation & 30 & Verified in the isotropic, normal-incidence, isothermal limit against published 2023/2024 dipolar-gradient band structures (agreement $\le x\%$), plus closed-form analytic dispersion and \STR\ eight internal consistency checks; all benchmark evidence is reported in the manuscript itself. \\
6 & Result 1 & 25 & Band gaps become direction-dependent: the first complete gap opens along $\Gamma$--X at $\theta=0^\circ$ and migrates to $\Gamma$--M as $\theta\to 90^\circ$. \\
7 & Result 2 & 22 & The $(\theta,\mathrm{AR})$ design map gives gap-width tunability of $y\%$; at $\mathrm{AR}=1$ the response is exactly rotationally invariant (verified). \\
8 & Result 3 & 20 & Iso-frequency contours and group-velocity fields deviate from the phase direction by up to $z^\circ$, enabling wave steering. \\
9 & Micro-inertia & 18 & \STR\ Without $\ell_{\mathrm i}$ the phase velocity is unbounded as $k\to\infty$ (proved in Appendix~A, confirmed numerically); the transition is quantified. \\
10 & Implication & 22 & Orientation is a manufacturing-controllable band-gap design variable for directional vibration isolation. \\
\bottomrule
\end{tabularx}
}

\subsection{Highlights \bp{Elsevier mandatory: 3--5 bullets, max 85 characters each}}

\begin{itemize}[label=\textbullet]
\item \texttt{Bloch C1 FE for 2D anisotropic strain-gradient elasticity} \hfill (54 chars)
\item \texttt{Ellipsoidal microstructure gives two band-gap design knobs} \hfill (57)
\item \texttt{Band gaps become direction-dependent on microstructure angle} \hfill (59)
\item \texttt{Micro-inertia is required for bounded phase velocity} \hfill (51)
\item \texttt{Validated against 2023 and 2024 published gradient-elastic bands} \hfill (64)
\end{itemize}

\subsection{Keywords \bp{6--8; mix theory, method, application}}

\emph{Strain-gradient elasticity; Mindlin Form-II; Bloch--Floquet analysis; Elastic band gap;
Anisotropic microstructure; Micro-inertia; $C^1$ finite element; Bogner--Fox--Schmit element;
Phononic crystal; Iso-frequency contour}

\subsection{Nomenclature \bp{one table, immediately after the abstract}}

Group the symbols into five blocks so the table is scannable:

{\small
\begin{tabularx}{\textwidth}{@{}L{2.9cm} L{7.2cm} Y@{}}
\toprule
Block & Representative symbols & Note \\
\midrule
Kinematics & $u_i,\ \varepsilon_{ij},\ \eta_{ijk}$ & displacement, strain, strain gradient \\
Elastic / gradient moduli & $C_{ijkl},\ A_{ijklmn},\ a_1\ldots a_5$ & 2 + 5 isotropic gradient constants \\
Microstructure geometry & $l_1,l_2,l_3;\ \mathrm{AR};\ \theta;\ (A^{\mathsf T}\!A)_{ij}$ & semi-axes, aspect ratio, orientation, rotated length tensor \\
Inertia & $\rho,\ \ell_{\mathrm i}$ & mass density, micro-inertia length \\
Wave / Bloch & $\bm k,\ \omega,\ \bm a_1,\bm a_2,\ \bm b_1,\bm b_2,\ \bm v_{\mathrm g}$ & wavenumber, frequency, lattice and reciprocal vectors, group velocity \\
\bottomrule
\end{tabularx}
}

\newpage
%%==============================================================================
\section{SECTION-BY-SECTION BLUEPRINT}
%%==============================================================================

\bp{Page budgets assume single-spaced elsarticle review format, 11pt.
Total target: 30--34 pages including references and Appendices A--B.} \STR

%%------------------------------------------------------------------------------
\subsection{Section 1 --- Introduction \bp{$\approx$3.5 pp}}

\textbf{Purpose:} establish that (i) gradient elasticity is needed for wave filtering at
small scale, (ii) micro-inertia is needed for causality, and (iii) \STR\ \emph{to the best of our knowledge, orientation-dependence of gradient-elastic band gaps has not been studied}. End with the four contributions.

\begin{tabularx}{\textwidth}{@{}L{1.5cm} Y L{4.4cm}@{}}
\toprule
Sub-\S & Content & Must-cite \\
\midrule
1.1 & Microstructured media and the breakdown of classical elasticity when the wavelength approaches the microstructural length scale. Name the physical platforms: rolled/drawn metallic textures, fibre-reinforced composites, additively manufactured lattices. & Mindlin 1964, 1965; Mindlin \& Eshel 1968; Toupin 1962 \\
1.2 & Gradient-elasticity family and why Form II. \textbf{Key argument: gradient elasticity without micro-inertia predicts unbounded phase velocity} $\Rightarrow$ \STR\ loss of bounded wave speed at large $k$, i.e.\ physical inadmissibility (asymptotic proof in Appendix~A). This motivates $\ell_{\mathrm i}$ as a constitutive parameter, not a numerical device. & Papargyri-Beskou \& Beskos 2009; Askes \& Aifantis 2011; Polyzos \& Fotiadis 2012 \\
1.3 & Phononic crystals and elastic band gaps: Bragg scattering, Bloch theorem, irreducible Brillouin zone, complete vs.\ partial gaps. & Kushwaha et al.\ 1993; Yan et al.\ 2010 \\
1.4 & \textbf{State of the art in gradient-elastic band gaps --- \STR\ to the best of our knowledge all 1D or isotropic (Table~1).} Present as a small literature table (see Table~1). This is where the three validation anchors appear naturally. & Li, Wei \& Zhou 2016; \textbf{Li, Askes, Gitman, Krynkin \& Wei 2023}; \textbf{Li et al.\ 2024}; Hosseini \& Zhang 2021; Zheng et al. \\
1.5 & \textbf{The gap.} \STR\ Three explicit sentences, each tied to Table~1: (a) within the scope of our review, all published gradient-elastic band-gap results are 1D; (b) 2D studies use isotropic length scales only; (c) microstructural \emph{orientation} has, to the best of our knowledge, not been used as a band-gap design variable. & --- \\
1.6 & Contributions C1--C4 (verbatim from the Lock Sheet, Section~1.2). & --- \\
1.7 & Organisation paragraph. & --- \\
\bottomrule
\end{tabularx}

\textbf{Planned float:} \textbf{Table 1} --- literature positioning table (see Section~\ref{sec:tables}).
\textbf{Planned equations:} none.
\textcolor{warn}{\textbf{Reviewer will check:}} does Table~1 make the gap self-evident without you having to assert it?

%%------------------------------------------------------------------------------
\subsection{Section 2 --- Anisotropic strain-gradient continuum with micro-inertia \bp{$\approx$4.5 pp}}

\textbf{Purpose:} the complete constitutive model, derived from scratch. This is the
theoretical core and must stand alone.

\begin{tabularx}{\textwidth}{@{}L{1.5cm} Y L{3.4cm}@{}}
\toprule
Sub-\S & Content & Equations \\
\midrule
2.1 & \textbf{Ellipsoidal nonlocal averaging domain.} Define the averaging neighbourhood $\mathcal E$ as an ellipsoid with semi-axes $(l_1,l_2,l_3)$. Define the characteristic-length tensor as the second moment of $\mathcal E$; show $(A^{\mathsf T}\!A)_{mn}=l_m^2\delta_{mn}$ in principal axes. Emphasise the geometric (fabric-tensor) interpretation. & (1)--(4) \\
2.2 & \textbf{Orientation.} Rotation $\bm R(\theta)$ about $z$; transformed tensor $(A^{\mathsf T}\!A)_{\mathrm{rot}}=\bm R^{\mathsf T}\mathrm{diag}(l_1^2,l_2^2,l_3^2)\bm R$; write out all components; identify the off-diagonal $(1,2)$ term and the mixed $\partial^2/\partial x_1\partial x_2$ contribution it generates. Give the in-plane closed form $g_{22}^2(\theta)=(l_1^2\sin^2\theta+l_2^2\cos^2\theta)/5$. & (5)--(9) \\
2.3 & \textbf{Kinematics.} Small strain $\varepsilon_{ij}=\tfrac12(u_{i,j}+u_{j,i})$; strain gradient $\eta_{ijk}=\varepsilon_{ij,k}$; count independent components in 2D plane strain. & (10)--(12) \\
2.4 & \textbf{Mindlin Form-II constitutive relations.} Cauchy stress $\sigma_{ij}=C_{ijkl}\varepsilon_{kl}$; double stress $\tau_{ijk}=A_{ijklmn}\eta_{lmn}$; isotropic sixth-order gradient modulus written with the five constants $a_1\ldots a_5$; plane-strain reduction. & (13)--(18) \\
2.5 & \textbf{Micro-inertia.} Kinetic energy density $T=\tfrac12\rho\,\dot u_i\dot u_i+\tfrac12\rho\,\ell_{\mathrm i}^2\,\dot u_{i,j}\dot u_{i,j}$. Derive the inertial term $\rho(\ddot u_i-\ell_{\mathrm i}^2\ddot u_{i,jj})$. \textbf{State explicitly that this is what bounds the phase velocity;} \STR\ the high-$k$ asymptotic proof is deferred to Appendix~A. & (19)--(21) \\
2.6 & \textbf{Strain energy density and positive definiteness.} Full $W$; sufficient conditions on $(a_1\ldots a_5)$ and on $(A^{\mathsf T}\!A)_{\mathrm{rot}}$; the anisotropic gradient stiffness in the form $K_g=\tfrac{1}{10}\sum_{i}\sum_{j}(A^{\mathsf T}\!A)_{\mathrm{rot},ij}\,\bm B_{,i}^{\mathsf T}\bm C\,\bm B_{,j}$. & (22)--(26) \\
2.7 & \textbf{Specialisations (the limit ladder).} Tabulate four reductions: (i) $l\to0,\ \ell_{\mathrm i}\to0$ = classical elasticity; (ii) $\mathrm{AR}=1$ = isotropic gradient elasticity; (iii) $\ell_{\mathrm i}\to0$ = gradient elasticity without micro-inertia (\STR\ unbounded wave speed, Appendix~A); (iv) $\theta\to\theta+90^\circ$ with $l_1\leftrightarrow l_2$ = symmetry. & (27)--(30) \\
2.8 & \textbf{Equations of motion and boundary conditions.} Strong form $\sigma_{ij,j}-\tau_{ijk,jk}=\rho(\ddot u_i-\ell_{\mathrm i}^2\ddot u_{i,jj})$; surface double traction $R_i=n_jn_k\tau_{ijk}$; the four boundary quantities per direction (classical traction, double traction, displacement, normal derivative). & (31)--(35) \\
\bottomrule
\end{tabularx}

\textbf{Planned float:} Figure~1 (ellipsoidal neighbourhood + rotated tensor schematic).
\textbf{Planned equations:} (1)--(35).
\textcolor{warn}{\textbf{Reviewer will check:}} positive definiteness under \emph{rotation}. Show that
rotation preserves the eigenvalues $\{l_1^2,l_2^2,l_3^2\}$ of $(A^{\mathsf T}\!A)$, hence
positive definiteness is orientation-independent --- a one-line but valuable result.

%%------------------------------------------------------------------------------
\subsection{Section 3 --- Bloch--Floquet formulation \bp{$\approx$2.5 pp}}

\textbf{Purpose:} convert the infinite-medium wave problem into a unit-cell eigenproblem,
and define every observable that will be reported.

\begin{tabularx}{\textwidth}{@{}L{1.5cm} Y L{3.4cm}@{}}
\toprule
Sub-\S & Content & Equations \\
\midrule
3.1 & Unit cell $\Omega_{\mathrm{cell}}=[0,L]\times[0,L]$; lattice vectors $\bm a_1,\bm a_2$; reciprocal vectors $\bm b_1,\bm b_2$; first Brillouin zone. & (36)--(38) \\
3.2 & \textbf{Bloch theorem for a $C^1$ gradient-elastic medium.} Time-harmonic ansatz; $\bm u(\bm x+\bm a_\alpha)=\bm u(\bm x)\,e^{\mathrm i\bm k\cdot\bm a_\alpha}$. \textbf{Then the critical step:} because the weak form contains first derivatives as independent DOFs, the \emph{same} phase factor must be imposed on $\nabla\bm u$ and on $\partial^2\bm u/\partial x\partial y$. Derive this from the Bloch condition rather than asserting it. & (39)--(43) \\
3.3 & Irreducible Brillouin zone for the square lattice: $\Gamma(0,0)\to X(\pi/L,0)\to M(\pi/L,\pi/L)\to\Gamma$. Discretisation of the path and the number of $k$-points per segment. & (44) \\
3.4 & \textbf{Non-dimensionalisation.} $\bar k=kL/\pi$, $\bar\omega=\omega/\omega_0$, $\omega_0=\sqrt{\mu/(\rho L^2)}$, $\bar l_m=l_m/L$, $\bar\ell_{\mathrm i}=\ell_{\mathrm i}/L$, $\bar v=\omega/k\,/\sqrt{\mu/\rho}$. \bp{This is what makes the anchor comparison work --- see Section~\ref{sec:validation}.} & (45)--(47) \\
3.5 & \textbf{Observables, defined precisely.} Band $\bar\omega_n(\bm k)$; band gap $\Delta\bar\omega_g=\bar\omega_{n+1}^{\min}-\bar\omega_n^{\max}$ over the whole IBZ (complete) or along one path (partial); gap mid-frequency and aspect ratio; \STR\ normalised gap width $\Delta\bar\omega/\bar\omega_{\mathrm{mid}}$ and orientation sensitivity $S_\theta=\max|\partial\Delta\bar\omega/\partial\theta|/\max\Delta\bar\omega$; \STR\ gaps exist only with periodic contrast (Case C) --- a homogeneous cell shows dispersion but no Bragg gap; phase velocity; group velocity $\bm v_{\mathrm g}=\nabla_{\bm k}\omega$; \STR\ time-averaged energy flux derived from the energy balance in Appendix~B (classical + double-stress terms, dimensional check), with $\bm v_{\mathrm g}=\langle\bm S\rangle/\langle W+T\rangle$ verified numerically in \S7.3. & (48)--(52) \\
\bottomrule
\end{tabularx}

\textbf{Planned float:} Figure~2 (unit cell, lattice, IBZ with $\Gamma$-X-M-$\Gamma$ path).
\textcolor{warn}{\textbf{Reviewer will check:}} the energy flux expression. In gradient elasticity the
Poynting-type vector has an extra double-stress term; omitting it is a common error.
Derive it from the energy balance and state the extra term explicitly.

%%------------------------------------------------------------------------------
\subsection{Section 4 --- $C^1$ finite element discretisation \bp{$\approx$3.5 pp}}

\textbf{Purpose:} the numerical method, with the Bloch-phase handling as the headline.

\begin{tabularx}{\textwidth}{@{}L{1.5cm} Y L{3.4cm}@{}}
\toprule
Sub-\S & Content & Equations \\
\midrule
4.1 & \textbf{Bogner--Fox--Schmit bicubic Hermite rectangle.} Tensor product of 1D cubic Hermite; 4 DOF per node per displacement component $\{u,\,u_{,x},\,u_{,y},\,u_{,xy}\}$; 2 components $\times$ 4 nodes $\times$ 4 DOF $=$ \textbf{32 DOF per cell}; $C^1$ inter-element continuity, conforming in $H^2(\Omega)$. Shape functions and the $\bm B$ / $\bm B_{,i}$ matrices. & (53)--(56) \\
4.2 & Element classical stiffness $\bm K^{\mathrm c}$ and gradient stiffness $\bm K^{\mathrm g}(\theta,\mathrm{AR})$; Gauss--Legendre order required (the integrand contains second derivatives of cubics $\Rightarrow$ up to quartic; state the rule used and justify it). & (57)--(59) \\
4.3 & \textbf{Element mass and gradient-inertia matrices} $\bm M_0$ and $\bm M^{\mathrm g}$; total $\bm M=\bm M_0+\ell_{\mathrm i}^2\bm M^{\mathrm g}$. & (60)--(61) \\
4.4 & \textbf{Bloch-periodic master--slave elimination.} Partition cell DOFs into interior / edge / corner. For each periodic pair, $\bm d_{\mathrm{slave}}=\bm T(\bm k)\,\bm d_{\mathrm{master}}$ with $\bm T$ diagonal complex phase. Assemble the transformation and form the reduced matrices. State the resulting reduced DOF count. & (62)--(65) \\
4.5 & \textbf{Reduced Hermitian eigenproblem} $[\bar{\bm K}(\bm k)-\bar\omega^2\bar{\bm M}(\bm k)]\bar{\bm d}=\bm 0$; prove $\bar{\bm K}(\bm k)=\bar{\bm K}(-\bm k)^{\mathsf H}$; solver choice and the number of lowest bands extracted. & (66)--(68) \\
4.6 & \textbf{Scaling, conditioning and mode tracking.} Jacobi (diagonal) scaling; report the condition number $\kappa$ vs $(\theta,\mathrm{AR})$; branch tracking along the IBZ path by eigenvector modal-assurance-criterion (MAC) continuation, with step halving on branch jump. & (69)--(71) \\
\bottomrule
\end{tabularx}

\textbf{Planned floats:} Figure~3 (BFS element DOF layout + Bloch phase on derivative DOFs).
\textbf{Planned equation count:} (53)--(71).

\begin{center}
\begin{minipage}{0.94\textwidth}
\rowcolors{2}{softamber}{softamber}
\begin{tabularx}{\textwidth}{@{}L{4.0cm} Y@{}}
\multicolumn{2}{@{}l@{}}{\textbf{\color{warn}Known numerical pitfalls --- design these tests in \emph{before} coding}}\\
Non-Hermitian $\bar{\bm K}$ & Wrong sign of the phase on the derivative DOFs. Test $\|\bar{\bm K}(\bm k)-\bar{\bm K}(-\bm k)^{\mathsf H}\|/\|\bar{\bm K}\|<10^{-12}$ on every run. \\
Spurious high-frequency branches & A $C^1$ bicubic cell carries more branches than the physical count. Report only the lowest $N$ and document the filter. \\
Condition number blow-up & Large $\mathrm{AR}$ or large $\bar l$ couples derivative DOFs across directions. Jacobi scaling; plot $\kappa$. \\
Branch crossing / mis-ordering & Use MAC-based continuation, not raw eigenvalue sorting. \\
Unbounded phase velocity & Only occurs if $\ell_{\mathrm i}=0$ --- this is a \emph{result}, not a bug (Section~7.4). \STR\ Appendix~A gives the analytic reason; Fig.~13(b) the numerical evidence. \\
Interface conditions (Case C) & Gradient elasticity needs \emph{four} conditions per interface (displacement, normal derivative, traction, double traction). Missing one silently shifts the gap edges. \\
\end{tabularx}
\end{minipage}
\end{center}

%%------------------------------------------------------------------------------
\subsection{Section 5 --- Verification \bp{$\approx$4.5 pp}}\label{sec:validation}

\textbf{Purpose:} make it impossible for a reviewer to say ``the method is not validated''. \STR\ All validation evidence (configurations, parameters, reproduced curves, error tables) appears in the manuscript itself.
See Section~\ref{sec:valmatrix} for the full \COR\ five-layer validation and verification framework; the structure inside the paper is:

\begin{xltabular}{\textwidth}{@{}L{1.5cm} Y L{3.4cm}@{}}
\toprule
Sub-\S & Content & Anchor \\
\midrule
5.1 & Verification strategy: state the \COR\ \textbf{five-layer validation and verification framework} up front: \COR\ \textbf{Layers 1--4 = external/published or analytical validation; Layer 5 = internal numerical verification + mesh convergence.} \STR\ \textbf{Terminology, stated explicitly:} Layers 1--4 \emph{validate the physical formulation} (published benchmarks, closed-form analytics, independent published method); Layer 5 \emph{verifies the numerical implementation} (eight-test suite + mesh convergence). The two are never conflated, and the internal tests are never described as published validation. One paragraph + one summary table. & --- \\
5.2 & \textbf{Layer 1 --- classical elasticity limit.} Set $\bar l\to0,\ \bar\ell_{\mathrm i}\to0$, bilayer unit cell, normal incidence. Reproduce the published classical dispersion branches and Bragg gaps. & \textbf{Li et al.\ 2024, Sci.\ Rep.\ 14:24035, Fig.~2(a)}; cross-ref.\ Zheng \& Wei 2009 \\
5.3 & \textbf{Layer 2 --- gradient-elasticity published benchmarks (the decisive test).} \COR\ Three benchmark configurations \emph{within this single layer} (2a, 2b, 2c --- not three separate studies):
(a) isotropic gradient elasticity with flexoelectricity suppressed, $\bar l=10^{-5}$, $\bar\ell_{\mathrm i}=2\times10^{-5}$, AlN/BaTiO$_3$ bilayer $\Rightarrow$ reproduce their Fig.~2(b);
(b) dipolar gradient elasticity, isothermal, Pb/brass bilayer, $a_1=10^{-5}$\,m $\Rightarrow$ reproduce their Fig.~4(c);
(c) band gaps for the classical / gradient / gradient-thermoelastic cases (first two only) $\Rightarrow$ reproduce their Fig.~3 (optional configuration).
\textbf{Method difference is the point: they used a transfer matrix, we use a $C^1$ Bloch FE.} \STR\ \textbf{Independent reproduction:} reference values are recomputed from the anchors' own equations, parameters and interface conditions with an independently written transfer-matrix code; digitised curves are used only to draw overlays, never to compute errors. & \textbf{Li et al.\ 2024, Fig.~2(b)} and \textbf{Li, Askes, Gitman, Krynkin \& Wei 2023, WRAM 36(4):5715--5735, Fig.~4(c)} \\
5.4 & \textbf{Layer 3 --- closed-form analytical validation.} Homogeneous medium (Case H): compare the computed longitudinal and shear dispersion against the analytic gradient-elastic relations for an infinite medium and an axial bar; \COR\ include the high-wavenumber asymptotic limit of Appendix~A. Machine-precision target. & Papargyri-Beskou \& Beskos 2009, Eqs.\ (22)--(28); Li, Wei \& Zhou 2016 \\
5.5 & \textbf{Layer 4 --- independent-method cross-check.} Homogeneous limit ($\mathrm{AR}=1$, no gradation) compared with a dynamic-stiffness / Floquet--Bloch computation on a lattice. & Mishra, Kumar \& Sharma 2026, Acta Mech.\ 237:3951--3982 \\
5.6 & \textbf{Layer 5 --- internal numerical verification and mesh convergence.} \COR\ The eight-test suite (5a--5h) and the mesh-convergence/resolution-floor study (\S5.7) form \emph{one} verification layer, not separate studies. \STR\ Eight automated tests, reported as a table: Hermiticity; BZ periodicity $\omega(\bm k+\bm G)=\omega(\bm k)$; rotational invariance at $\mathrm{AR}=1$; long-wave acoustic slope $\omega\to c|\bm k|$; $\theta\!\to\!\theta\!+\!90^\circ$ symmetry with $l_1\leftrightarrow l_2$; positive definiteness of $\bar{\bm K},\bar{\bm M}$; bounded phase velocity for $\bar\ell_{\mathrm i}>0$; \STR\ energy-flux/group-velocity identity $\langle\bm S\rangle=(\langle W\rangle+\langle T\rangle)\bm v_{\mathrm g}$ vs $\nabla_{\bm k}\omega$ (5h). & self \\
5.7 & \textbf{Layer 5 (continued) --- mesh convergence and resolution floor.} \COR\ Bands at $4{\times}4,8{\times}8,16{\times}16,32{\times}32$ BFS cells; report the converged first-gap edge and the \emph{numerical resolution floor} $\varepsilon_\Delta$ --- the smallest gap width distinguishable from noise. \textbf{This is what allows weak anisotropy effects to be claimed as real.} \STR\ All four prescribed refinement levels are computed and reported with their eigenvalues, relative errors and step-to-step changes. The observed convergence rate is obtained by least squares from the levels whose relative error exceeds the numerical resolution of the measurement: a level participates in the fit if and only if its relative error is greater than the pre-declared factor $F = 3$ times the larger of that level's measured reproducibility and the declared numerical zero $10^{-15}$: $e_i > F \cdot \max(s_i, 10^{-15})$, $F = 3$ --- the reproducibility $s_i$ being measured from two pre-registered start vectors at the frozen solver configuration (the protocol, $F$ and the declared zero are fixed before the calculation is run and recorded with it). Levels that do not meet this criterion are reported as resolution-limited data, with their relative errors, their measured reproducibility and the resulting ratio, and are excluded from the fit; the identity of any excluded level is stated in the text, the table and the figure caption. The 95\% CI of the fit is reported; a theoretical order is claimed only if justified by conforming-subspace eigenvalue theory (Babu\v{s}ka--Osborn), otherwise the observed rate is reported as-is, together with the number of fitted levels. If fewer than three levels satisfy the criterion, no rate is reported. & self \\
\bottomrule
\end{xltabular}

\textbf{Planned floats:} Figure~4 (anchor comparison, 2 panels), Figure~5 (mesh convergence + resolution floor), Table~3 (anchor error table), Table~4 (internal consistency suite).

\begin{center}
\begin{minipage}{0.94\textwidth}
\rowcolors{2}{softred}{softred}
\begin{tabularx}{\textwidth}{@{}L{3.6cm} Y@{}}
\multicolumn{2}{@{}l@{}}{\textbf{\color{warn}HARD GATE G3 --- non-negotiable}}\\
Criterion & \COR\ Layers 1, 2a and 2b must each be validated by one of the two admissible evidence routes of Section~13: \textbf{(i) quantitative validation} --- band-gap edges and the first four branch frequencies agree with \emph{source numerical values} to \textbf{$\le 2\%$} relative (target $\le 0.5\%$ in the classical limit); or \textbf{(ii) graphical validation} --- where no machine-readable source values exist, the published curve is reproduced from the source-reported parameters and compared as an explicitly labelled overlay, with \emph{no} numerical percentage asserted. Numerical source data always outrank graphical evidence; a benchmark met by neither route is \textbf{not validated} and fails this gate. A benchmark that is
\textbf{not validated} solely because the published source is insufficient, and that satisfies all of the
conditions of Section~13 A3, is retained in that state and does not of itself fail this gate. \\
Evidence in manuscript & \STR\ Overlay/side-by-side dispersion plots (Fig.~4); tabulated first four branch frequencies at stated $\bar k$; band-gap edge comparison; relative error per quantity; maximum relative error; explicit PASS/FAIL per benchmark --- all in the main text, never supplementary-only. \\
If it fails & \textbf{Do not submit.} Debug using the troubleshooting table in \texttt{ANCHOR\_DATA\_SHEET.md}. The usual culprits, in order: (i) phase sign on derivative DOFs; (ii) micro-inertia sign --- the published equation of motion is $\sigma_{zz,z}-\mu_{zzz,zz}=\rho(\ddot u_z-\ell_1^2\ddot u_{z,zz})$, note the \emph{minus}; (iii) non-dimensionalisation mismatch --- use their exact $\omega_0$; (iv) interface condition count. \\
\end{tabularx}
\end{minipage}
\end{center}

\begin{center}
\begin{minipage}{0.94\textwidth}
\rowcolors{2}{softgreen}{softgreen}
\begin{tabularx}{\textwidth}{@{}L{3.4cm} Y@{}}
\multicolumn{2}{@{}l@{}}{\textbf{\color{ok}INDEPENDENT REPRODUCTION PROTOCOL} \ADDED}\\
Reference data & Recompute each anchor's dispersion from its published equations, parameters and interface conditions using an independently written transfer-matrix code (or published tables where available). Curve digitisation is permitted \emph{only} to draw overlay figures, never to compute error numbers. \\
Preserved exactly & Layer thicknesses and geometry; material constants; constitutive assumptions (isothermal dipolar gradient; flexoelectricity suppressed); the four interface conditions per boundary; the anchor's scaling, $\omega_0$ and non-dimensional axes. \\
Error definition & \textbf{Quantitative route only}, against \emph{source numerical values}: $e=|\bar\omega_{\mathrm{ours}}-\bar\omega_{\mathrm{ref}}|/\bar\omega_{\mathrm{ref}}$ per branch and per gap edge; report per-quantity and maximum values; PASS iff $\le2\%$. Where only the published \emph{graph} exists, no error is computed and none may be stated: the benchmark is recorded as \textbf{graphical validation} (Section~13), never as a percentage, and plotted similarity may not be converted into an error value. \\
\end{tabularx}
\end{minipage}
\end{center}

%%------------------------------------------------------------------------------
\subsection{Section 6 --- Band structure results \bp{$\approx$4.5 pp}}

\begin{tabularx}{\textwidth}{@{}L{1.5cm} Y L{3.4cm}@{}}
\toprule
Sub-\S & Content & Float \\
\midrule
6.1 & Baseline configuration: two unit-cell cases. \textbf{Case H} = homogeneous cell (dispersion only, used for analytic verification). \textbf{Case C} = matrix with a centred circular inclusion (impedance contrast $\Rightarrow$ Bragg gaps). Full parameter table. \textcolor{warn}{Note: a homogeneous medium does \emph{not} produce complete Bragg gaps --- this must be stated, not discovered mid-review.} \STR\ \textbf{Two mechanisms stated up front:} (i) homogeneous-medium microstructure-induced dispersion (Case H, no gaps); (ii) Bragg gaps from periodic impedance contrast (Case C). The manuscript never attributes a Bragg gap to a homogeneous medium. & Table~2 \\
6.2 & Case H results: microstructure-induced dispersion; comparison of the first four bands at $\mathrm{AR}=1$ vs $\mathrm{AR}=10$; where in the BZ the effect is largest. \STR\ Labelled ``dispersion'', never ``gap'', for Case H. & Fig.~6 \\
6.3 & Case C baseline band structure along $\Gamma$-X-M-$\Gamma$; identification of the first three complete and partial gaps; Bloch mode shapes at the gap edges. & Fig.~7 \\
6.4 & \textbf{Orientation sweep} at fixed $\mathrm{AR}=5$: $\theta\in\{0^\circ,15^\circ,\ldots,90^\circ\}$. \textbf{Headline result: the gap migrates from $\Gamma$--X to $\Gamma$--M.} & Fig.~8 \\
6.5 & \textbf{Aspect-ratio sweep} at fixed $\theta=45^\circ$: $\mathrm{AR}\in\{1,2,3,5,7,10\}$. Gap width vs $\mathrm{AR}$. & Fig.~9 \\
6.6 & \textbf{$(\theta,\mathrm{AR})$ band-gap design map} --- first complete gap width as a 3D response surface with contour projection; $7\times6=42$ computed points. & Fig.~10 \\
6.7 & Polar design map in $(X=\mathrm{AR}\cos\theta,\ Y=\mathrm{AR}\sin\theta)$ and the complete-vs-partial gap regime classification; the directional filtering criterion stated as an inequality. \STR\ Table~5 carries normalised gap width $\Delta\bar\omega/\bar\omega_{\mathrm{mid}}$ and $S_\theta$ columns. & Fig.~11, Table~5 \\
\bottomrule
\end{tabularx}

%%------------------------------------------------------------------------------
\subsection{Section 7 --- Anisotropic wave steering and physical admissibility \bp{$\approx$2.5 pp}}

\begin{tabularx}{\textwidth}{@{}L{1.5cm} Y L{3.4cm}@{}}
\toprule
Sub-\S & Content & Float \\
\midrule
7.1 & \textbf{Iso-frequency contours} at three frequencies $\times$ three orientations; show the transition from circular ($\mathrm{AR}=1$) to strongly non-convex ($\mathrm{AR}=10$, $\theta=45^\circ$). & Fig.~12 \\
7.2 & \textbf{Group velocity vs phase direction.} Compute $\bm v_{\mathrm g}=\nabla_{\bm k}\omega$ by central differences on the band surface; plot the deviation angle $\delta=\angle\bm v_{\mathrm g}-\angle\bm k$ vs propagation direction for each $\theta$. \STR\ Report $\delta_{\max}$ and a steering figure-of-merit (angular shift of the stop-band centre vs $\theta$) as numbers in the abstract and conclusions. & Fig.~12 \\
7.3 & Time-averaged energy flux, including the double-stress contribution derived in \S3.5 and Appendix~B; energy partition between classical and gradient parts as a function of frequency. \STR\ Numerically verify $\bm v_{\mathrm g}=\langle\bm S\rangle/\langle W+T\rangle$ against central-difference $\nabla_{\bm k}\omega$ (test 5h). & Fig.~13 \\
7.4 & \textbf{Micro-inertia and physical admissibility.} \STR\ Phase velocity vs $k$ for $\bar\ell_{\mathrm i}=0$ and $\bar\ell_{\mathrm i}>0$; demonstrate the unbounded branch in the former \emph{alongside the Appendix~A asymptotic proof} (analytical + numerical support, never numerics alone); phrase as loss of bounded wave speed, not as a numerics-only causality claim. Tie back to \S1.2 and to the published 2024 finding. & Fig.~13 \\
\bottomrule
\end{tabularx}

%%------------------------------------------------------------------------------
\subsection{Section 8 --- Discussion \bp{$\approx$1.5 pp}}
\begin{itemize}
\item \textbf{8.1 Physical mechanism.} Why orientation controls the gap: the boundary-normal component of the rotated length tensor penalises transverse strain gradients, so the effective stop-band follows the direction in which the gradient energy is largest. Give the one-line scaling argument, not a hand-wave.
\item \textbf{8.2 Manufacturing route and design implications.} Rolled/drawn textures and additively manufactured lattices have an \emph{inherently oriented} microstructure, so $\theta$ is set by the process, not by an extra part. Directional vibration-isolation mounts and wave-filtering panels. \bp{tie to the project's Haryana industrial-relevance section here --- one sentence only, no overclaiming}
\item \textbf{8.3 Limitations and scope.} Linear, isothermal, non-piezoelectric, 2D, perfect interfaces, lossless (no damping $\Rightarrow$ real $\omega$); graded and dissipative extensions and the inclusion of memory-type thermal coupling are out of scope. \bp{this paragraph is what makes the paper honest and also positions the follow-up work without promising it here}
\end{itemize}

%%------------------------------------------------------------------------------
\subsection{Section 9 --- Conclusions \bp{$\approx$1 pp}}
Seven numbered conclusions, one per contribution plus three quantitative findings.
Every conclusion must carry a number that appears in the results. No vague statements.

%%------------------------------------------------------------------------------
\subsection{Appendices A and B \bp{manuscript appendices, separate equation numbering}}

\ADDED\ These appendices exist so that the two physics claims a reviewer will attack
first --- unbounded phase velocity without micro-inertia, and the higher-order energy
flux --- rest on derivations, not on plots. Their equations are numbered (A.$\cdot$) and
(B.$\cdot$), so the main-text register (1)--(71) is unchanged.

\begin{tabularx}{\textwidth}{@{}L{1.5cm} Y L{3.4cm}@{}}
\toprule
App. & Content & Equations \\
\midrule
A & \textbf{High-wavenumber asymptotics.} Leading-order dispersion as $\bar k\to\infty$ with and without micro-inertia; proof that $\bar v\to\infty$ when $\bar\ell_{\mathrm i}=0$ and remains bounded when $\bar\ell_{\mathrm i}>0$; dimensional-consistency check; the admissibility statement used in \S1.2, \S2.7 and \S7.4. & (A.1)--(A.6) \\
B & \textbf{Energy-flux derivation.} Energy-balance identity; classical and double-stress (higher-order) flux terms; time-averaging; dimensional consistency; the identity $\bm v_{\mathrm g}=\langle\bm S\rangle/\langle W+T\rangle$ used in \S7.3 and internal test 5h. & (B.1)--(B.8) \\
\bottomrule
\end{tabularx}

\newpage
%%==============================================================================
\section{EQUATION REGISTER \bp{planned numbering --- keep it stable while writing}}
%%==============================================================================

\begin{longtable}{@{}>{\bfseries}p{1.5cm} L{7.4cm} L{5.4cm}@{}}
\toprule
No. & Object & Appears in \\
\midrule
\endfirsthead
\toprule
No. & Object & Appears in \\
\midrule
\endhead
(1) & Ellipsoidal averaging domain $\mathcal E$ and its indicator & \S2.1 \\
(2) & Nonlocal (averaged) field definition & \S2.1 \\
(3) & Second-moment characteristic-length tensor, principal axes & \S2.1 \\
(4) & $(A^{\mathsf T}\!A)_{mn}=l_m^2\delta_{mn}$ & \S2.1 \\
(5) & Rotation matrix $\bm R(\theta)$ about $z$ & \S2.2 \\
(6) & $(A^{\mathsf T}\!A)_{\mathrm{rot}}=\bm R^{\mathsf T}\mathrm{diag}(l_1^2,l_2^2,l_3^2)\bm R$ & \S2.2 \\
(7) & Expanded component form incl.\ off-diagonal $(1,2)$ term & \S2.2 \\
(8) & Axial weighting $(A^{\mathsf T}\!A)_{\mathrm{rot},11}=l_1^2\cos^2\theta+l_2^2\sin^2\theta$ & \S2.2 \\
(9) & In-plane closed form $g_{22}^2(\theta)=(l_1^2\sin^2\theta+l_2^2\cos^2\theta)/5$ & \S2.2 \\
(10)--(12) & $\varepsilon_{ij}$, $\eta_{ijk}$, independent-component count & \S2.3 \\
(13)--(15) & $\sigma_{ij}=C_{ijkl}\varepsilon_{kl}$; $\tau_{ijk}=A_{ijklmn}\eta_{lmn}$; Voigt form & \S2.4 \\
(16)--(18) & Isotropic sixth-order gradient modulus ($a_1\ldots a_5$); plane-strain reduction & \S2.4 \\
(19) & Kinetic energy density with micro-inertia & \S2.5 \\
(20) & Inertial operator $\rho(\ddot u_i-\ell_{\mathrm i}^2\ddot u_{i,jj})$ & \S2.5 \\
(21) & Hamilton's principle including the gradient-inertia term & \S2.5 \\
(22)--(24) & Strain energy density $W$; positive-definiteness conditions & \S2.6 \\
(25) & Rotation preserves eigenvalues $\{l_1^2,l_2^2,l_3^2\}$ & \S2.6 \\
(26) & $K_g=\tfrac{1}{10}\sum_i\sum_j(A^{\mathsf T}\!A)_{\mathrm{rot},ij}\bm B_{,i}^{\mathsf T}\bm C\,\bm B_{,j}$ & \S2.6 \\
(27)--(30) & Four specialisations (classical / isotropic / no micro-inertia / symmetry) & \S2.7 \\
(31) & Strong form of the equations of motion & \S2.8 \\
(32)--(35) & Double traction $R_i$, four BC quantities per direction, higher-order BCs & \S2.8 \\
(36)--(38) & Lattice vectors, reciprocal vectors, first BZ & \S3.1 \\
(39) & Time-harmonic Bloch ansatz & \S3.2 \\
(40) & $\bm u(\bm x+\bm a_\alpha)=\bm u(\bm x)e^{\mathrm i\bm k\cdot\bm a_\alpha}$ & \S3.2 \\
(41)--(43) & \textbf{Phase factor on $\nabla\bm u$ and $\partial^2\bm u/\partial x\partial y$} & \S3.2 \\
(44) & IBZ path $\Gamma$-X-M-$\Gamma$ and $k$-point sampling & \S3.3 \\
(45)--(47) & Non-dimensionalisation ($\bar k,\bar\omega,\bar l,\bar\ell_{\mathrm i},\bar v$) & \S3.4 \\
(48) & Band $\bar\omega_n(\bm k)$ & \S3.5 \\
(49) & Complete vs partial band gap, width, mid-frequency, aspect ratio, \STR\ normalised width, $S_\theta$ & \S3.5 \\
(50) & Phase velocity & \S3.5 \\
(51) & Group velocity $\bm v_{\mathrm g}=\nabla_{\bm k}\omega$ & \S3.5 \\
(52) & \STR\ Energy-balance identity; classical + double-stress flux; time-averaged form; $\bm v_{\mathrm g}=\langle\bm S\rangle/\langle W+T\rangle$ & \S3.5, App.~B \\
(53)--(56) & BFS shape functions; $\bm B$; $\bm B_{,i}$; DOF layout (32/cell) & \S4.1 \\
(57)--(59) & $\bm K^{\mathrm c}$, $\bm K^{\mathrm g}(\theta,\mathrm{AR})$, integration rule & \S4.2 \\
(60)--(61) & $\bm M_0$, $\bm M^{\mathrm g}$, $\bm M=\bm M_0+\ell_{\mathrm i}^2\bm M^{\mathrm g}$ & \S4.3 \\
(62)--(65) & Bloch transformation $\bm T(\bm k)$; master--slave elimination; reduced matrices & \S4.4 \\
(66)--(68) & Reduced Hermitian eigenproblem; $\bar{\bm K}(\bm k)=\bar{\bm K}(-\bm k)^{\mathsf H}$ & \S4.5 \\
(69)--(71) & Jacobi scaling; condition number; MAC continuation criterion & \S4.6 \\
(A.1)--(A.3) & \ADDED\ Leading-order high-$\bar k$ asymptotics; unbounded branch for $\bar\ell_{\mathrm i}=0$ & App.~A \\
(A.4)--(A.6) & \ADDED\ Bounded limit $\bar\ell_{\mathrm i}>0$; dimensional check; admissibility statement & App.~A \\
(B.1)--(B.4) & \ADDED\ Energy balance; classical flux; double-stress flux; time-averaged form & App.~B \\
(B.5)--(B.8) & \ADDED\ Dimensional check; $\langle\bm S\rangle/(\langle W\rangle+\langle T\rangle)=\bm v_{\mathrm g}$; numerical verification protocol & App.~B \\
\bottomrule
\end{longtable}

\textbf{Total: 71 numbered equations in the main text} \STR\ + (A.1)--(A.6) and (B.1)--(B.8) in Appendices A--B \ADDED.
\bp{That is high but appropriate for IJMS/AMM, which reward complete derivations. If the editor asks to shorten, move (16)--(18) and (62)--(65) to an Appendix --- never delete them.}

\newpage
%%==============================================================================
\section{FIGURE REGISTER \bp{13 figures}}
\label{sec:figures}
%%==============================================================================

\begin{xltabular}{\textwidth}{@{}>{\bfseries}p{1.05cm} L{4.15cm} Y L{2.15cm}@{}}
\toprule
Fig. & Title & Panels & Source of data \\
\midrule
\endfirsthead
\toprule
Fig. & Title & Panels & Source of data \\
\midrule
\endhead
1 & Ellipsoidal nonlocal averaging neighbourhood and the rotated characteristic-length tensor & (a) ellipsoid with semi-axes $l_1,l_2$ at $\theta=0^\circ,45^\circ,90^\circ$; (b) polar plot of $(A^{\mathsf T}\!A)_{\mathrm{rot},11}$ vs $\theta$; (c) unit cell with inclusion & schematic \\
2 & Periodic medium, lattice vectors and the irreducible Brillouin zone & (a) lattice; (b) BZ with $\Gamma$-X-M-$\Gamma$ & schematic \\
3 & BFS element and the Bloch phase on $C^1$ degrees of freedom & (a) 32-DOF layout; (b) master--slave pairs with phase factors on value and derivative DOFs & schematic \\
4 & \textbf{Verification against published results} & (a) classical limit vs 2024 Fig.~2(a); (b) gradient-elastic limit vs 2024 Fig.~2(b); (c) isothermal dipolar-gradient limit vs 2023 Fig.~4(c) & \textbf{anchors A, B} \\
5 & Mesh convergence and numerical resolution floor & (a) first gap edge vs mesh; (b) \STR\ observed convergence rate fitted from the levels that exceed the pre-declared resolution criterion (\S5.7; 95\% CI) --- the number of fitted levels and any excluded level are stated --- no theoretical order claimed unless justified; (c) resolution floor $\varepsilon_\Delta$ & computed \\
6 & Case H: microstructure-induced dispersion in a homogeneous medium & (a) bands along $\Gamma$-X-M-$\Gamma$; (b) $\mathrm{AR}=1$ vs $\mathrm{AR}=10$ overlay & computed \\
7 & Case C: baseline band structure & (a) bands; (b)--(d) Bloch mode shapes at the first three gap edges & computed \\
8 & \textbf{Orientation sweep at $\mathrm{AR}=5$} & (a) stacked bands for 7 values of $\theta$; (b) first gap edge along $\Gamma$-X vs $\theta$; (c) same along $\Gamma$-M & computed \\
9 & \textbf{Aspect-ratio sweep at $\theta=45^\circ$} & (a) bands for 6 values of AR; (b) gap width vs AR & computed \\
10 & \textbf{$(\theta,\mathrm{AR})$ band-gap design map} & 3D response surface + bottom contour projection; 42 computed points marked & computed \\
11 & Polar design map and gap-regime classification & (a) gap width in $(X,Y)$; (b) complete-vs-partial regime boundaries & computed \\
12 & \textbf{Iso-frequency contours and wave steering} & (a)--(c) contours at 3 frequencies for 3 orientations; (d) deviation angle $\delta$ vs propagation direction & computed \\
13 & Energy flux and micro-inertia admissibility & (a) classical vs gradient energy partition vs frequency; (b) phase velocity vs $k$ for $\bar\ell_{\mathrm i}=0$ and $\bar\ell_{\mathrm i}>0$ & computed \\
\bottomrule
\end{xltabular}

\textbf{Figure style rules} (so the set looks like one paper):
\begin{itemize}
\item Vector PDF/EPS only. No raster screenshots.
\item One consistent colour set; colour-blind safe; every curve also distinguished by marker/line style.
\item Font in figures $\ge$ footnote size of the journal (8pt minimum at final width).
\item Every panel labelled (a), (b), \ldots\ and every panel referenced in the text.
\item Axes: non-dimensional quantities with bars ($\bar\omega$, $\bar k$) unless a dimensional comparison against an anchor requires units --- then state the units.
\item Each figure $\le$ half a page; Figs.~4, 10, 12 may be full width.
\end{itemize}

%%==============================================================================
\section{TABLE REGISTER \bp{6 tables}}
\label{sec:tables}
%%==============================================================================

\begin{xltabular}{\textwidth}{@{}>{\bfseries}p{1.05cm} L{4.0cm} Y@{}}
\toprule
Table & Title & Columns / content \\
\midrule
\endfirsthead
\toprule
Table & Title & Columns / content \\
\midrule
\endhead
1 & \textbf{Literature positioning} (in the Introduction) & Ref.\ $\cdot$ year $\cdot$ dimension (1D/2D/3D) $\cdot$ theory $\cdot$ micro-inertia? $\cdot$ anisotropic length scale? $\cdot$ orientation studied? $\cdot$ method $\cdot$ \STR\ \textbf{reported wave/band-gap quantity} (dispersion relation / partial gap / complete gap / iso-frequency contour / group velocity / directional steering). \textbf{Our row is the only one with all four ``yes''} \STR\ and the only one populating the last three quantity columns. \\
2 & Material, geometric and microstructural parameters & Case H and Case C: $E$, $\nu$, $\rho$, $l_1$, $l_2$, $\ell_{\mathrm i}$, $\mathrm{AR}$, $\theta$, $L$, inclusion radius and contrast. \STR\ With a \textbf{provenance tag} on every value: [C] taken from a cited published source, [A] analytically defined/derived, [S] explicitly assumed with justification; validation rows carry the anchor's exact values [C-anchor]. No untagged number may reach the results. \\
3 & \textbf{Anchor comparison --- quantitative} & \STR\ Band $\cdot$ BZ location $\cdot$ first four branch frequencies $\cdot$ gap edges $\cdot$ published $\bar\omega$ $\cdot$ computed $\bar\omega$ $\cdot$ relative error \% per quantity $\cdot$ maximum relative error $\cdot$ PASS/FAIL. Rows for anchor A Fig.~4(c) and anchor B Figs.~2(a),(b). \\
4 & \textbf{Internal consistency suite} & Test $\cdot$ quantity $\cdot$ expected $\cdot$ computed $\cdot$ pass/fail. \STR\ Eight rows (incl.\ energy-flux/group-velocity identity 5h). \COR\ (Layer 5.) \\
5 & Band-gap summary & Gap order $\cdot$ type (complete/partial) $\cdot$ lower edge $\cdot$ upper edge $\cdot$ width $\cdot$ mid-frequency $\cdot$ aspect ratio, for each $(\theta,\mathrm{AR})$ combination. \STR\ + normalised gap width $\Delta\bar\omega/\bar\omega_{\mathrm{mid}}$ and orientation sensitivity $S_\theta$. \\
6 & Convergence and resolution & Mesh $\cdot$ DOF $\cdot$ first gap edge $\cdot$ relative change $\cdot$ $\kappa$ $\cdot$ wall-clock. Plus the benchmarked $\varepsilon_\Delta$. \STR\ + observed convergence rate (least-squares fit, 95\% CI); a theoretical order appears only if justified. \COR\ (Layer 5.) \\
\bottomrule
\end{xltabular}

\newpage
%%==============================================================================
\section{VALIDATION MATRIX}
\label{sec:valmatrix}
%%==============================================================================

\COR\ \textbf{Five-layer validation and verification framework.} Layers 1--4 = external/published or analytical validation of the physical formulation; Layer 5 = internal numerical verification (eight-test suite 5a--5h) \emph{plus} mesh convergence and the resolution floor (row 5i). Subcases 2a/2b/2c are three benchmark configurations \emph{within} Layer 2, not three separate validation studies. The internal tests of Layer 5 are never described as published validation.

\rowcolors{2}{softblue}{white}
\begin{xltabular}{\textwidth}{@{}c L{5.0cm} L{2.6cm} c Y@{}}
\toprule
Layer & Test & Anchor & Year & Status \\
\midrule
1 & Classical elasticity limit ($\bar l=\bar\ell_{\mathrm i}=0$), 1D bilayer, normal incidence --- reproduce published dispersion and Bragg gaps & Li, Li, Guo et al., \emph{Sci.\ Rep.} \textbf{14}:24035, Fig.~2(a) & \textbf{2024} & \textcolor{warn}{\textbf{GATE}} \\
2a & Gradient elasticity, flexoelectricity suppressed ($\bar l=10^{-5}$, $\bar\ell_{\mathrm i}=2\times10^{-5}$, $f=0$), AlN/BaTiO$_3$ bilayer --- reproduce published Fig.~2(b) & Li et al., \emph{Sci.\ Rep.} \textbf{14}:24035, Fig.~2(b) & \textbf{2024} & \textcolor{warn}{\textbf{GATE}} \\
2b & Dipolar gradient elasticity, \textbf{isothermal} (thermoelastic coupling ignored), Pb/brass bilayer, $a_1=10^{-5}$\,m --- reproduce published Fig.~4(c) & Li, Askes, Gitman, Krynkin \& Wei, \emph{Waves Random Complex Media} \textbf{36}(4):5715--5735 & \textbf{2023} & \textcolor{warn}{\textbf{GATE}} \\
2c & Band gaps for the three cases: classical / gradient / gradient-thermoelastic (first two only) & same, Fig.~3 & \textbf{2023} & optional \\
3 & Closed-form analytic dispersion for infinite medium and axial bar, gradient elasticity with micro-elastic $g^2$ and micro-inertia $h^2$ & Papargyri-Beskou \& Beskos, \emph{IJSS} \textbf{46}:2151--2159, Eqs.\ (22)--(28) & 2009 & analytic \\
3b & Purely elastic 1D dipolar-gradient phononic crystal & Li, Wei \& Zhou, \emph{Acta Mech.} \textbf{227}:1005--1023 & 2016 & analytic \\
4 & Homogeneous limit cross-checked against a different numerical method (dynamic stiffness + Wittrick--Williams) & Mishra, Kumar \& Sharma, \emph{Acta Mech.} \textbf{237}:3951--3982 & \textbf{2026} & optional \\
5a & Hermiticity $\|\bar{\bm K}(\bm k)-\bar{\bm K}(-\bm k)^{\mathsf H}\|/\|\bar{\bm K}\|<10^{-12}$ & self & --- & mandatory \\
5b & BZ periodicity $\omega(\bm k+\bm G)=\omega(\bm k)$ & self & --- & mandatory \\
5c & Rotational invariance at $\mathrm{AR}=1$ & self & --- & mandatory \\
5d & Long-wave acoustic slope $\omega\to c|\bm k|$, relative error $<10^{-6}$ & self & --- & mandatory \\
5e & $\theta\to\theta+90^\circ$ symmetry with $l_1\leftrightarrow l_2$ & self & --- & mandatory \\
5f & Positive definiteness of $\bar{\bm K}$ and $\bar{\bm M}$ & self & --- & mandatory \\
5g & Bounded phase velocity for $\bar\ell_{\mathrm i}>0$ (and unbounded for $=0$) & self & --- & mandatory \\
5h & Energy-flux/group-velocity identity $\langle\bm S\rangle=(\langle W\rangle+\langle T\rangle)\bm v_{\mathrm g}$ vs $\nabla_{\bm k}\omega$ \ADDED & self & --- & mandatory \\
5i \COR & Mesh convergence $4^2\!\to\!32^2$ + resolution floor $\varepsilon_\Delta$ --- part of Layer 5, not a sixth layer & self & --- & mandatory \\
\bottomrule
\end{xltabular}



\medskip
\textbf{\color{ok}Why this satisfies the mandatory condition.} Three separate published
papers from \textbf{2023}, \textbf{2024} and \textbf{2026} are used as external references,
two of them as hard submission gates. Layer 2b is the decisive one: the published
configuration is \emph{dipolar gradient elasticity with thermoelastic coupling ignored},
which is algebraically the same model as this paper in the 1D normal-incidence limit ---
but computed by a transfer matrix rather than by a $C^1$ Bloch finite element.
Agreement therefore validates the formulation \emph{and} the discretisation simultaneously.
\STR\ Terminology used throughout: external Layers 1--4 = \emph{validation of the physical formulation}; \COR\ Layer 5 (eight-test suite + mesh convergence) = \emph{verification of the numerical implementation}; the internal tests are never described as published validation.

\newpage
%%==============================================================================
\section{REFERENCE PLAN \bp{target 62--72 references}}
%%==============================================================================

\begin{tabularx}{\textwidth}{@{}L{5.0cm} >{\centering\arraybackslash}p{1.4cm} Y@{}}
\toprule
Category & Count & Notes \\
\midrule
Validation anchors (2023, 2024, 2026) & 3 & \textbf{Must be cited in the Introduction \emph{and} in Section 5.} \\
Gradient / strain-gradient elasticity foundations & 8 & Mindlin 1964, 1965; Mindlin \& Eshel 1968; Toupin 1962; Tiersten 1962; Germain 1973; Askes \& Aifantis 2011 \\
Wave dispersion in gradient elasticity & 6 & Papargyri-Beskou \& Beskos 2009; Polyzos \& Fotiadis 2012; Georgiadis and co-workers; Li, Wei \& Zhou 2016 \\
Phononic crystals / band gaps general & 8 & Kushwaha 1993; Yan 2010; review + 5 recent \\
Gradient-elastic band gaps (the direct prior art) & 6 & incl.\ Hosseini \& Zhang 2021 (meshless collocation), Zheng et al. \\
Bloch--Floquet FE / dynamic stiffness methods & 5 & incl.\ Mishra et al.\ 2026; wave-FE and scaled-boundary variants \\
$C^1$ finite elements & 4 & Bogner--Fox--Schmit 1965/1968; Hermite; higher-order-continuity implementations \\
Eigenvalue/convergence theory \ADDED & 2 & Babu\v{s}ka \& Osborn 1991 (conforming eigenvalue approximation); a standard FE-eigenvalue convergence reference --- cited only if a convergence order is claimed. \\
Micro-inertia / gradient inertia & 3 & incl.\ the 2024 anchor's explicit statement \\
Elastic metamaterials / vibration isolation applications & 6 & 2022--2026 preferred \\
Anisotropy in phononic band gaps & 4 & Zhan \& Wei 2010, 2014 (anisotropic band gaps --- classical theory) \\
\midrule
\textbf{Total} & \textbf{55--74} & Aim for $\ge$40\% published 2021 or later. \\
\bottomrule
\end{tabularx}

\textbf{\color{warn}Citation hygiene:} every reference must be actually cited in the text at
least once; no orphan entries. Anchor A and Anchor B must appear in \S1.4 (as prior art),
in Table~1 (positioning), and in \S5.2--5.3 (as validation). Three appearances each.

\medskip
\textbf{Must-cite shortlist (full strings ready for BibTeX in \texttt{ANCHOR\_DATA\_SHEET.md}):}
\begin{enumerate}[label=\arabic*.]
\item Li, Askes, Gitman, Krynkin \& Wei (2023), \emph{Waves Random Complex Media} 36(4):5715--5735. \textbf{[2023 anchor]}
\item Li, Li, Guo et al.\ (2024), \emph{Sci.\ Rep.} 14:24035. \textbf{[2024 anchor]}
\item Mishra, Kumar \& Sharma (2026), \emph{Acta Mech.} 237:3951--3982. \textbf{[2026 anchor]}
\item Li, Wei \& Zhou (2016), \emph{Acta Mech.} 227:1005--1023.
\item Papargyri-Beskou \& Beskos (2009), \emph{IJSS} 46:2151--2159.
\item Zheng \& Wei (2009), \emph{Int.\ J.\ Miner.\ Metall.\ Mater.} 16(5):608--614.
\item Mindlin (1964), \emph{ARMA} 16:51--78; Mindlin \& Eshel (1968), \emph{IJSS} 4:109--124.
\item Zhan \& Wei (2010), \emph{Acta Mech.\ Solida Sin.} 23:182--188 --- anisotropy and band gaps, classical theory; cite to show that anisotropy matters \emph{even classically}, and that gradient-elastic anisotropy is unstudied.
\item Hosseini \& Zhang (2021), \emph{Int.\ J.\ Mech.\ Sci.} 209:106711 --- meshless band-structure method.
\item Askes \& Aifantis (2011), \emph{IJSS} 48:3482--3490 --- gradient elasticity overview incl.\ FE implementation.
\end{enumerate}

%%==============================================================================
\section{WRITING CONVENTIONS \bp{fix these before writing a single paragraph}}
%%==============================================================================

\begin{tabularx}{\textwidth}{@{}L{4.2cm} Y@{}}
\toprule
Item & Convention \\
\midrule
Index notation & Latin $i,j,k,l,m,n\in\{1,2,3\}$; Einstein summation throughout; comma $=$ partial derivative. \\
Vectors / tensors & Bold lowercase $\bm u,\bm k,\bm v_{\mathrm g}$; bold uppercase $\bm R,\bm C,\bm K,\bm M,\bm B,\bm T$. \\
Imaginary unit & $\mathrm i$ (roman), never $i$. \\
Micro-inertia symbol & $\ell_{\mathrm i}$ --- \textbf{not} $l_1$, to avoid collision with the characteristic-length semi-axes $l_1,l_2,l_3$. \textcolor{warn}{The 2024 anchor uses $l_1$ for micro-inertia; when comparing, state the mapping explicitly in a footnote.} \\
Orientation angle & $\theta$ (this paper). The anchors use $\alpha$; the comparison table must say so. \\
Aspect ratio & $\mathrm{AR}=l_1/l_2$ (roman, not a product of $A$ and $R$). \\
Non-dimensional & Over-bar: $\bar\omega,\bar k,\bar l,\bar\ell_{\mathrm i},\bar v$. Dimensional quantities carry units. \\
Units & SI throughout. Report both dimensional and non-dimensional where an anchor comparison needs it. \\
Significant figures & 4 in tables of results, 3 in running text. Never more than the convergence justifies. \\
Time dependence & $e^{-\mathrm i\omega t}$ --- state once in \S3.2 and never change it. \\
Numbers at sentence start & Spell out. \\
Figure references & Always ``Fig.~7(b)'', never ``the figure below''. \\
Passive voice & Preferred; ``we'' permitted in \S1.6 and \S9 only. \\
Tense & Present for established results and for our own findings; past for what the anchors did. \\
\STR\ Validation vs verification & \COR\ \emph{Validation} = Layers 1--4: agreement with external published/analytical/independent-method references (the formulation describes reality); \emph{verification} = Layer 5: internal eight-test suite + mesh convergence (the equations are solved correctly). The words are never used interchangeably; internal tests are never called published validation. \\
\STR\ Novelty claims & Absolute claims (``never'', ``nobody'', ``has not been studied'') are forbidden unless tied to Table~1; use ``to the best of our knowledge'' / ``within the scope of our review''. \\
\STR\ Parameter provenance & Every numerical parameter carries [C] cited / [A] analytically defined / [S] assumed-with-justification. No untagged number reaches the results. \\
\STR\ Homogeneous vs periodic & ``Band gap'' is used only for periodic contrast (Case C); the homogeneous Case H exhibits microstructure-induced dispersion only. \\
\bottomrule
\end{tabularx}

\newpage
%%==============================================================================
\section{PRODUCTION PLAN AND CHECKLISTS}
%%==============================================================================

\subsection{Fourteen-week schedule with four gates}

\begin{longtable}{@{}>{\bfseries}p{1.5cm} L{8.2cm} L{4.6cm}@{}}
\toprule
Week & Task & Output \\
\midrule
\endfirsthead
\toprule
Week & Task & Output \\
\midrule
\endhead
1 & \textbf{DONE.} Anchors downloaded, parameters extracted, data sheet written. & \texttt{ANCHOR\_DATA\_SHEET.md} \\
2 & Section 2 derivation: \S2.1--\S2.4 (ellipsoidal tensor, rotation, Form-II constitutive). & Draft \S2, Eqs.\ (1)--(18) \\
3 & \S2.5--\S2.8 (micro-inertia, energy, positive definiteness, EOM, BCs). Symbolic degree check with exact rational arithmetic. & \textbf{GATE G1} --- symbolic checks pass \\
4 & BFS element: shape functions, $\bm B$, $\bm B_{,i}$, $\bm K^{\mathrm c}$, $\bm K^{\mathrm g}$. Patch test. & Element module \\
5 & $\bm M_0$, $\bm M^{\mathrm g}$, Bloch transformation, master--slave elimination. & Bloch module \\
6 & Hermiticity test, BZ sweep, MAC branch tracking. Case H runs. & \textbf{GATE G2} --- all 8 automated internal tests pass \\
7 & \STR\ Re-implement each anchor's transfer-matrix dispersion independently from its published equations and parameters; digitise figures \emph{only} for overlays; implement their exact non-dimensionalisation. & Anchor comparison harness \\
8 & Layers 1, 2a, 2b, 2c executed. Error table produced for every quantity comparable against source numerical values; a labelled graphical-validation record produced for every benchmark without such values. & \textbf{GATE G3} --- \COR\ each of the three gates met either quantitatively at $\le2\%$ or, where source numerical values do not exist, by the graphical route of Section~13 with its classification recorded; the route used is stated per benchmark \\
9 & \COR\ Layers 3, 4 and 5 (analytic, cross-method, internal suite + mesh convergence, $\varepsilon_\Delta$). & Section 5 complete \\
10 & Case C production runs: baseline + $\theta$ sweep (7) + AR sweep (6) $=$ 42 analyses. & Figs.~7--9 \\
11 & Design maps + polar map + gap-regime classification. & Figs.~10--11, Table~5 \\
12 & Iso-frequency contours, group velocity, energy flux, micro-inertia study. & Figs.~12--13 \\
13 & Write \S1, \S3, \S4, \S6--\S9. Assemble. & Full draft \\
14 & Abstract, highlights, cover letter, figure polish, reference audit, self-plagiarism re-check. & \textbf{GATE G4} --- submission-ready \\
\bottomrule
\end{longtable}

\subsection{Pre-submission checklist}

\textbf{Manuscript}
\begin{itemize}
\item \chk Title, all authors, affiliations, emails, ORCID, corresponding author marked
\item \chk \textbf{HSRF sanction number 14/66-2025 SPD/HSHEC/HSRF, Session 2025--2026} in the funding statement
\item \chk Abstract 230--260 words, no citations, no undefined abbreviations
\item \chk Highlights 3--5, each $\le$85 characters
\item \chk Keywords 6--8
\item \chk Nomenclature table present
\item \chk ``Implemented in MATLAB R20xxb / Python 3.x'' stated; code + data availability with a real URL/DOI
\item \chk CRediT statement, competing-interest declaration
\item \chk Limitations paragraph in \S8.3
\end{itemize}

\textbf{Validation (the mandatory condition)}
\begin{itemize}
\item \chk \textbf{2023 anchor} (WRAM 36(4):5715--5735, Fig.~4(c)) reproduced, error tabulated
\item \chk \textbf{2024 anchor} (Sci.\ Rep.\ 14:24035, Figs.~2(a) and 2(b)) reproduced, error tabulated
\item \chk \textbf{2026 anchor} (Acta Mech.\ 237:3951--3982) homogeneous-limit cross-check
\item \chk Closed-form analytic check to machine precision
\item \chk \STR\ All eight internal consistency tests reported as a table
\item \chk \STR\ Benchmark configurations, exact parameters, model assumptions, non-dimensionalisation, reproduced curves and error tables present in the manuscript (not supplementary-only)
\item \chk \STR\ First four branch frequencies and gap edges tabulated with per-quantity and maximum relative error and explicit PASS/FAIL
\item \chk \STR\ Every parameter carries a provenance tag [C]/[A]/[S]
\item \chk \STR\ Observed convergence rate reported; no unjustified theoretical order claimed
\item \chk \STR\ Micro-inertia claim supported by Appendix~A asymptotics \emph{and} Fig.~13(b) numerics
\item \chk \STR\ Homogeneous-vs-periodic mechanism distinction stated in \S3.5 and \S6.1
\item \chk Mesh convergence and resolution floor $\varepsilon_\Delta$ reported
\item \chk Every anchor cited three times (Introduction, Table 1, Section 5)
\end{itemize}

\textbf{Uniqueness}
\begin{itemize}
\item \chk Table~1 literature positioning shows the four ``yes'' columns unique to this paper
\item \chk Micro-inertia and the $C^1$ Bloch-phase-on-derivatives listed as explicit contributions
\item \chk No dependence on any unpublished manuscript of ours
\item \chk Text-overlap self-check against the eight existing PDFs re-run, $<$2\%
\end{itemize}

\textbf{Production quality}
\begin{itemize}
\item \chk All 13 figures vector, consistent style, 8pt minimum, every panel referenced
\item \chk All 6 tables have units and significant figures within convergence
\item \chk Every equation numbered is referenced at least once
\item \chk No orphan references; $\ge$40\% of references from 2021 or later
\item \chk Language pass (British or American English, consistent)
\end{itemize}

\subsection{Journal-specific deltas}
\label{sec:submission}

\begin{tabularx}{\textwidth}{@{}L{4.6cm} Y@{}}
\toprule
If submitting to & Change \\
\midrule
Int.\ J.\ Mechanical Sciences (primary) & Nothing. Full derivations welcome; keep all 71 equations. Emphasise C3 (the $C^1$ Bloch element) and C4 (design map) in the cover letter. \\
J.\ Sound and Vibration & Shorten \S2 by moving (16)--(18) to an Appendix. Strengthen \S7 (wave steering) and \S8.2 (vibration isolation). Use title Alternate A. \\
Composite Structures & Add one paragraph in \S8.2 on architected/lattice media and AM. Use title Alternate B. \\
Appl.\ Math.\ Modelling & Keep all derivations; expand \S5 (verification) by half a page; state the numerical-analysis contributions more formally. \\
Waves Random Complex Media (fallback) & Shorten to $\approx$24 pp; drop \S7.3; use title Alternate A. Strategic value: anchors A and C are in this journal. \\
\bottomrule
\end{tabularx}

%%------------------------------------------------------------------------------
\subsection{Publication-critical requirements \bp{hard submission gate}}
\label{sec:pcr}

\ADDED\ The manuscript \textbf{cannot be submitted} until every line below is true and
evidenced in the computation log. This subsection is the single enforcement point for
the publication-readiness review recorded in Section~\ref{sec:gcr}.

\begin{enumerate}[label=\textbf{PCR\arabic*.}]
\item All mandatory published benchmarks (Layers 1, 2a, 2b) are validated by an admissible
evidence route (Section~13): \textbf{quantitative validation} with maximum relative error
$\le2\%$ wherever source numerical values exist (the $0.5\%$ target applies only to the
classical limit), or \textbf{graphical validation} where they do not --- explicitly labelled,
with no numerical percentage asserted. A benchmark that is neither is \textbf{not validated}
and this item fails, \emph{unless} all of the conditions of Section~13 A3 hold for that benchmark, in
which case the benchmark is retained in the same \textsc{Not validated} state, reported with the
recorded source-side reason, and does not of itself fail this item; a graphical result must never be
reported or described as quantitative.
\item Benchmark configurations, exact parameters, model assumptions, non-dimensionalisation,
reproduced dispersion curves, numerical comparisons and relative errors all appear in the
\emph{main manuscript}, not only in supplementary or internal files.
\item Analytical checks pass: closed-form dispersion (Layer 3), Appendix~A high-$k$
asymptotics, Appendix~B energy-flux identity.
\item \COR\ The eight-test internal consistency suite (Layer~5, together with PCR5) passes and is tabulated (Table~4).
\item Mesh convergence is demonstrated with the \emph{observed} rate (least-squares fit over the refinement levels admissible under the pre-declared resolution criterion of \S5.7, 95\% CI) and the resolution floor $\varepsilon_\Delta$ is reported (Fig.~5, Table~6); every level excluded from the fit is reported with its relative error, measured reproducibility and ratio. \COR\ (Layer~5.)
\item Every material, geometric and microstructural parameter carries a provenance tag
[C]/[A]/[S] (Table~2); no untagged number reaches the results.
\item Every major physical claim is quantitatively supported: orientation-controlled gaps
(normalised gap widths, $S_\theta$, Table~5), wave steering ($\delta_{\max}$ and a steering
figure-of-merit, Fig.~12), micro-inertia necessity (Appendix~A proof \emph{and} Fig.~13(b)).
\item Novelty statements are hedged per the writing conventions, and the research gap is
objectively visible in Table~1 including the reported-quantity column.
\end{enumerate}

\medskip
\textbf{Gate owner:} the PI signs G4 only after PCR1--PCR8 are checked against the
computation log. A failed PCR blocks submission exactly as G3 does.

%%==============================================================================
\section{GAP CLOSURE REPORT \bp{revision 1.0 $\to$ 1.1, 22 September 2026}}
\label{sec:gcr}

Scope of the review: twelve publication-readiness requirements for Q1, IF\,$>$\,5
submission (validation gate; independent reproduction; validation evidence; micro-inertia
claim; novelty wording; literature positioning; parameter provenance; homogeneous vs
periodic; quantitative interpretation; convergence; energy flux; final checklist).
\textbf{Nothing was rewritten or restructured;} every change below carries \ADDED\ or
\STR\ in the body of this document.

\subsection*{(a) Already sufficient in v1.0 --- kept unchanged}
\begin{itemize}
\item Hard gate G3 with an explicit do-not-submit rule and the $\le2\%$ tolerance.
\item Five-layer validation ladder \COR\ (naming consolidated to ``five-layer validation and verification framework'' in revision 1.2, Section~12) with three published anchors (2023, 2024, 2026), GATE markers, and the method-difference argument (transfer matrix vs $C^1$ Bloch FE).
\item Anchor overlays (Fig.~4) and an anchor error table (Table~3) already planned as
manuscript floats; mesh convergence + resolution floor; closed-form analytic layer.
\item Parameter table (Table~2) already demanded provenance; the homogeneous-no-Bragg-gap
warning already existed in \S6.1; the energy-flux reviewer-check note already existed in \S3.
\item Self-containment rule (no citation of the unpublished FEM\_1--5 manuscripts).
\end{itemize}

\subsection*{(b) Missing in v1.0 --- now \ADDED}
\begin{itemize}
\item \textbf{Independent Reproduction Protocol} box in \S5: reference values recomputed
from the anchors' own equations; digitisation restricted to overlay drawing; exact-preservation
list; explicit error definition and PASS rule.
\item \textbf{Validation-evidence row} inside the G3 box: overlays, first four branch
frequencies, gap-edge comparison, per-quantity and maximum relative error, PASS/FAIL, and
the in-manuscript (not supplementary) rule.
\item \textbf{Appendix~A} (high-$k$ asymptotics, (A.1)--(A.6)) and \textbf{Appendix~B}
(energy-flux derivation, (B.1)--(B.8)); new blueprint subsection 3.10 and register rows.
\item \textbf{Eighth internal test 5h} (energy-flux/group-velocity identity) in \S5.6, the
validation matrix and Table~4.
\item \textbf{Quantitative steering metrics}: normalised gap width and $S_\theta$ (\S3.5,
Table~5); $\delta_{\max}$ and a steering figure-of-merit (\S7.2).
\item \textbf{Provenance tag scheme} [C]/[A]/[S] (Table~2, writing conventions, checklist, PCR6).
\item \textbf{Reported-quantity column} in Table~1 (dispersion relation / partial gap /
complete gap / IFC / group velocity / steering).
\item \textbf{Subsection 10.4} Publication-Critical Requirements (PCR1--PCR8).
\item \textbf{Terminology rule} validation (formulation) vs verification (implementation)
in \S5.1, the matrix note and the conventions.
\item This Gap Closure Report.
\end{itemize}

\subsection*{(c) Insufficient in v1.0 --- now \STR}
\begin{itemize}
\item \textbf{Novelty wording:} D4, C3, \S1 purpose, \S1.2, \S1.4, \S1.5, abstract S2 ---
absolute claims (``nobody'', ``never'', ``has not been documented'') replaced by ``to the
best of our knowledge / within the scope of Table~1''.
\item \textbf{Micro-inertia claim:} C2, abstract S9, \S1.2, \S2.5, \S2.7, \S7.4 ---
``non-causal'' replaced by loss of bounded wave speed, supported analytically (Appendix~A)
\emph{and} numerically (Fig.~13(b)); causality language never from numerics alone.
\item \textbf{Validation-in-manuscript rule:} D5, \S5 purpose, abstract S5, checklist, PCR2.
\item \textbf{Convergence:} \S5.7, Fig.~5(b), Table~6 --- observed rate fitted with 95\% CI;
a theoretical order only if justified (Babu\v{s}ka--Osborn added to the reference plan).
\item \textbf{Table column sets:} Table~2 (provenance tags), Table~3 (first four branches,
gap edges, per-quantity and maximum error, PASS/FAIL), Table~5 (normalised width, $S_\theta$),
Table~6 (observed rate).
\item \textbf{Homogeneous vs periodic:} \S3.5, \S6.1, \S6.2 and a conventions row.
\item \textbf{Energy flux:} \S3.5 Eq.~(52) object, \S7.3, equation register.
\item \textbf{Counts corrected/updated:} internal tests $7\to8$; equation total now states
71 main-text + appendix equations; title-page planned length and equation count corrected
(v1.0 said $\approx$58 equations and 29--32 pp against the register's 71; now 71 + Appendices
A--B and 30--34 pp); reference total 53--72 $\to$ 55--74.
\end{itemize}

\subsection*{(d) Insertion map --- requirement $\to$ location}
{\small
\begin{tabularx}{\textwidth}{@{}c L{6.2cm} Y@{}}
\toprule
Req. & Action & Inserted at \\\midrule
1 & \STR\ + \ADDED & D5; \S5 purpose; \S5.1 terminology; G3 evidence row; checklist; PCR1--2 \\
2 & \ADDED\ + \STR & Reproduction Protocol box (\S5); \S5.3; Week~7 \\
3 & \STR\ + \ADDED & G3 evidence row; Table~3; Fig.~4; checklist; PCR2 \\
4 & \STR\ + \ADDED & C2; abstract S9; \S1.2; \S2.5; \S2.7; \S7.4; Appendix~A; pitfall box; PCR7 \\
5 & \STR & D4; C3; \S1 purpose, 1.2, 1.4, 1.5; abstract S2; conventions; PCR8 \\
6 & \STR & Table~1 register row (reported-quantity column) \\
7 & \STR\ + \ADDED & Table~2 row; conventions row; checklist; PCR6 \\
8 & \STR & \S3.5; \S6.1; \S6.2; conventions row \\
9 & \STR & \S3.5 Eq.~(49); \S6.7 / Table~5; \S7.2; PCR7 \\
10 & \STR & \S5.7; Fig.~5(b); Table~6; reference plan; PCR5 \\
11 & \STR\ + \ADDED & \S3.5 Eq.~(52); \S7.3; Appendix~B; test 5h; matrix row 5h \\
12 & \ADDED & Subsection 10.4 (PCR1--PCR8) \\
\bottomrule
\end{tabularx}}

\medskip
\textbf{Honesty statement.} No requirement above is claimed as already satisfied
numerically or experimentally: this document remains a locked \emph{plan}. Satisfaction is
recorded only when the corresponding computation exists and is logged (Weeks 2--14, gates
G1--G4, PCR1--PCR8).

%%==============================================================================
\section{CHANGE MADE \bp{revision 1.1 $\to$ 1.2, 22 September 2026}}
\label{sec:change12}

\COR\ Only the corrections below were performed. No restructuring, no new physics, no
new validation study, no change to the planned length/figures/tables/equations, and no
weakening of any publication-critical requirement.

\begin{enumerate}[label=\textbf{M\arabic*.}]
\item \textbf{Layer numbering corrected.} ``Five-layer validation ladder'' (and the
implicitly six-row matrix) replaced everywhere by the \textbf{five-layer validation and
verification framework}: Layer~1 classical published benchmark; Layer~2 gradient-elasticity
published benchmarks (configurations 2a/2b/2c); Layer~3 closed-form analytical validation
(incl.\ the Appendix~A high-$k$ asymptotic limit); Layer~4 independent published-method
cross-check; Layer~5 internal numerical verification (eight-test suite 5a--5h) \emph{plus}
mesh convergence and resolution floor. Locations: D5, \S5 intro, \S5.1, \S5.3--\S5.7,
validation-matrix lead and rows, Week~9, writing conventions, table registers, PCR4--5,
one-page summary.
\item \textbf{Matrix row renumbered.} Former Layer~6 (mesh convergence, $\varepsilon_\Delta$)
is now row 5i inside Layer~5; \S5.7 relabelled ``Layer~5 (continued)''. All eight tests, the
four mesh levels, the observed-rate fit with 95\% CI and the resolution-floor criterion are
unchanged.
\item \textbf{Layer~2 subcases declared configurations.} \S5.3 now lists configurations
(a), (b) and (c) and states that they are three benchmark configurations within one layer,
not three studies; the matrix lead repeats this.
\item \textbf{G3/PCR1 untouched in substance.} The gate still mandates quantitative
reproduction of Layers 1, 2a and 2b at $\le2\%$ maximum relative error ($\le0.5\%$ target
in the classical limit), with all benchmark evidence in the main manuscript; the G3
criterion line now names Layers 1, 2a, 2b explicitly.
\item \textbf{Validation/verification split restated.} Layers 1--4 = external validation of
the physical formulation; Layer~5 = internal verification of the numerical implementation;
internal tests are never described as published validation (D5, \S5.1, matrix lead and
note, writing conventions).
\item \textbf{Workload unchanged.} No new published benchmark, experiment, 3D/thermal/
flexoelectric/nonlinear study, or additional independent numerical method; the eight tests
remain one suite; mesh convergence remains part of Layer~5.
\end{enumerate}

\medskip
\textbf{Consistency audit (performed after the edits):}
\begin{itemize}
\item \YES\ Exactly five validation/verification layers are defined (matrix rows 1, 2a--2c, 3/3b, 4, 5a--5i).
\item \YES\ Layer~2 contains configurations 2a, 2b, 2c.
\item \YES\ Layer~5 contains the eight internal tests (5a--5h) plus mesh convergence/resolution floor (5i).
\item \YES\ G3/PCR1 require the published benchmarks of Layers 1, 2a, 2b: quantitatively at $\le2\%$ ($0.5\%$ classical target) where source numerical values exist, or by the labelled graphical route of Section~13 where they do not (amendment A2, v1.5; thresholds unchanged); evidence-in-manuscript rule intact. \\
\item \YES\ No PCR1--PCR8 requirement removed or weakened; all \ADDED/\STR\ safeguards of v1.1 remain in force.
\item \YES\ No unnecessary new validation study introduced; planned length (30--34 pp), figures (13), tables (6) and equations (71 + appendices) unchanged.
\end{itemize}

%%==============================================================================

%%==============================================================================
\section{AMENDMENT A2 \bp{revision 1.4 $\to$ 1.5, 24 September 2026 --- external-validation evidence routes}}

\textbf{Authority and scope.} \COR\ Amendment A2 was directed by the Principal Investigator (P12R brief,
24 September 2026) so that a published \emph{graphical} benchmark can satisfy the external-validation
requirement of Layers 1, 2a and 2b where the source provides no machine-readable numerical values. A2
changes \emph{how} a published benchmark may be evidenced. It does \textbf{not} change any numerical
criterion, threshold, model, parameter, mesh or production result, it creates no new benchmark, and it
promotes no gate (A2.8). Itemisation, the exact delta and the immutability record:
\texttt{paper9/audit/BLUEPRINT\_V1\_5\_GRAPHICAL\_VALIDATION\_AMENDMENT.md}.

\subsection*{A2.1 --- The three evidence states (exhaustive and mutually exclusive)}
\begin{enumerate}[label=\textbf{(\alph*)}]
\item \textbf{Quantitative validation:} source numerical values (tables, released data or the authors'
values) exist, the existing quantitative comparison is performed, the $\le 2\%$ criterion
($\le 0.5\%$ target in the classical limit) is applied unchanged, and the prescribed error
$e=|\bar\omega_{\mathrm{ours}}-\bar\omega_{\mathrm{ref}}|/\bar\omega_{\mathrm{ref}}$ is reported per
quantity and as a maximum.
\item \textbf{Graphical validation:} no machine-readable source values exist, but the published source
provides a sufficiently detailed benchmark graph together with the material, geometric, normalisation
and boundary/interface information needed to reconstruct the case.
\item \textbf{Not validated:} the published graph or its accompanying information is insufficient to
reproduce the case reliably, or the source suffers an unresolved ambiguity. No claim may be made and the
benchmark fails G3/PCR1 --- \emph{unless} all of the conditions of Section~13 A3 hold for that benchmark,
in which case the benchmark is retained in this same state, reported as \textsc{Not validated} for the
recorded source-side reason, and does not of itself fail that item.
\end{enumerate}
Every mandatory benchmark is reported in exactly one of these states, and the state is stated
explicitly per benchmark.

\subsection*{A2.2 --- Requirements of the graphical route (all mandatory)}
\begin{enumerate}[label=\textbf{G\arabic*.}]
\item Reproduce the benchmark using the source-reported material, geometric, normalisation and
boundary/interface information; parameters are never adjusted to improve the visual match.
\item Compare the reproduced curve \emph{directly} against the published curve.
\item Preserve the source's figure identification and all relevant axes and normalisations.
\item Document the observable characteristics used for the comparison: branch identity,
frequency/wavenumber range, turning points and band-gap features, and any quantity independent of the
drawing.
\item Provide an overlay or equivalent figure where technically possible.
\item Classify the result explicitly as \textsc{Graphical validation} --- never as
\textsc{Quantitative validation}.
\end{enumerate}

\subsection*{A2.3 --- Evidence hierarchy}
\textbf{Numerical source values $>$ published graphical source with sufficient parameters $>$
insufficient source.} The lower route must never be presented, tabulated or described as stronger
evidence than the higher route. If source numerical values exist, the graphical route is not available
for that benchmark.

\subsection*{A2.4 --- Unresolved source ambiguity (B2-specific)}
For the strain-gradient benchmark of the 2024 anchor (Fig.~2(b)) the published definition of the
internal length scale is ambiguous ($l$ versus $\bar l$). That ambiguity must be resolved from
authoritative source information before B2 can be accepted under \emph{either} route. No interpretation
may be chosen silently, and B2 may not be reported as validated while the ambiguity stands.

\subsection*{A2.5 --- Prohibition of fabricated precision}
Graphical agreement must never be converted into a numerical error. No percentage --- including
``$<2\%$'', ``$\approx1.x\%$'' or any bound --- may be attached to a graphically validated benchmark
unless source numerical values exist, or unless a separately governed, explicitly declared digitisation
protocol has been approved (A2 does not approve one). Curve digitisation remains permitted \emph{only}
to draw overlays, never to compute error numbers.

\subsection*{A2.6 --- Machine-readable classification contract}
The three states are fixed identifiers, usable by tooling and by the manuscript evidence table:
\texttt{QUANTITATIVE\_VALIDATION}, \texttt{GRAPHICAL\_VALIDATION}, \texttt{NOT\_VALIDATED}. A record in
the graphical state carries \emph{no} numerical error field; a record without source numerical values may
not carry one. Reference implementation and synthetic tests:
\texttt{paper9/verification/suite/benchmark\_validation\_route.py},
\texttt{paper9/verification/suite/test\_p12r\_graphical\_validation\_route.py}.

\subsection*{A2.7 --- Decision table}
\begin{center}
\begin{tabularx}{\textwidth}{@{}L{6.2cm} Y@{}}
\toprule
Evidence situation & Route and required treatment \\
\midrule
Source numerical values available & \textbf{Quantitative route}: existing comparison, existing thresholds
($\le2\%$; $0.5\%$ classical), prescribed error reported \\
\midrule
Graph only, parameters sufficient & \textbf{Graphical route (A2.2)}: reproduce, overlay, document
observables, classify \textsc{Graphical validation}; no percentage \\
\midrule
Graph only, parameters ambiguous or insufficient, and all conditions of Section~13 A3 satisfied &
\textbf{Not validated}: no claim; reported with the recorded source-side reason; the G3/PCR1 item is not
failed by this benchmark alone \\
\midrule
Not attempted, or not reproduced for a reason other than source insufficiency & \textbf{Not validated}:
no claim; the G3/PCR1 item fails for that benchmark \\
\bottomrule
\end{tabularx}
\end{center}

\subsection*{A3 --- Source-insufficient benchmarks: conditions and consequence}
\COR\ A benchmark is \textsc{Not validated} for a \emph{source-side} reason when the published
source does not provide the information needed to reproduce the case, or suffers an ambiguity that
cannot be resolved from the published record. Such a benchmark does not of itself fail PCR1
item~1 or Gate~G3 \emph{only if} every condition below holds. The conditions are conjunctive: if
any one fails, the benchmark keeps the ordinary consequence of Section~13 A2.1(c) and the item
fails.
\begin{enumerate}[label=\textbf{A3.\arabic*.}, leftmargin=*]
\item \textbf{A validated benchmark exists.} At least one mandatory external benchmark
(Layers~1, 2a, 2b) is validated by an admissible route under Section~13 A2.7.
\item \textbf{Explicit, everywhere.} The \textsc{Not validated} state is stated explicitly, per
benchmark, in the main manuscript and in every active status-carrying artefact --- the machine
record, the anchor-error table, the status tables and the blocker matrices. No artefact may
imply that the benchmark is validated.
\item \textbf{No numerical claim.} No relative error, percentage, bound, digitised value or
inferred number is claimed for a source-insufficient benchmark; its
\texttt{quantitative\_error} stays null. A2.5 remains fully in force.
\item \textbf{Source-side attribution recorded.} The reason is recorded as source-side, with the
formulation-consistency evidence retained, and any unresolved ambiguity is recorded as
unresolved rather than silently resolved. A2.4 remains fully in force.
\item \textbf{Higher-tier route retained.} The prepared author-data request package is retained
as the higher-tier quantitative route. If author data arrive, the benchmark moves to the
quantitative route and must then meet the $\le2\%$ criterion (and the $0.5\%$ classical target
where it applies) exactly as before. Authorising A3 neither sends nor authorises sending
anything.
\end{enumerate}
\STR\ \textbf{No automatic promotion.} A3 re-scopes a consequence: it creates no evidence route,
admits no new kind of evidence, adds no benchmark and adds no validation requirement, and it
changes no threshold, no numerical result, no figure, no table and no scientific claim. PCR1 and
G3 are recorded as \textsc{Pass} and \textsc{Met} only by a separate reassessment against
A3.1--A3.5 recorded with its evidence; G4 remains a separate PI signature under Section~10.3 and
P13/submission remains a separate PI decision.

\subsection*{A2.8 --- Lines re-scoped by A2, historical records, and non-promotion}
\COR\ A2 re-scopes, in substance, the following v1.4 statements: the G3 hard-gate criterion row
(line~457) and its evidence row (line~458, whose per-quantity relative-error list applies to the extent
the route provides each item --- for a graphically validated benchmark the required evidence is the
overlay figure, the documented observables and the route classification); the error-definition row
(line~471); the Week-8 gate row (line~803); PCR1 (lines~886--888); and the consistency-audit bullet
(line~1046). Sections~11 and~12 are \emph{historical revision records} and are deliberately left
unamended; they describe earlier revisions and do not state current acceptance policy. The footer
revision banner keeps its historical wording; A2 appends its own note there.

\textbf{No automatic promotion.} The existence of A2 does not mark PCR1 or G3 as PASS/MET: PCR1 remains
\textbf{NOT PASS} and G3 remains \textbf{NOT MET} until an actual benchmark reproduction is performed and
audited under A2. A2 also leaves untouched the quantitative thresholds, the mathematical model, the
production results, the manuscript text (whose validation wording must later be re-tiered to match the
state actually achieved --- a separately authorised edit), and the prepared author-data request package,
which remains available as the higher-tier (quantitative) route.
\section*{\color{lockblue}ONE-PAGE SUMMARY}
%%==============================================================================
\addcontentsline{toc}{section}{One-page summary}

\rowcolors{2}{softgrey}{white}
\begin{xltabular}{\textwidth}{@{}L{3.3cm} Y@{}}
\textbf{Problem} & Two-dimensional Bloch wave propagation and band-gap formation in a strain-gradient elastic medium whose microstructure is an \emph{oriented ellipsoid}, including gradient (micro-) inertia. \\
\textbf{Governing equation} & $\sigma_{ij,j}-\tau_{ijk,jk}=\rho\left(\ddot u_i-\ell_{\mathrm i}^2\ddot u_{i,jj}\right)$ \\
\textbf{Anisotropy} & $(A^{\mathsf T}\!A)_{\mathrm{rot}}=\bm R^{\mathsf T}(\theta)\,\mathrm{diag}(l_1^2,l_2^2,l_3^2)\,\bm R(\theta)$, \quad $\mathrm{AR}=l_1/l_2$ \\
\textbf{Discretisation} & Bogner--Fox--Schmit $C^1$ rectangle, 32 DOF/cell, complex Bloch phase on value \emph{and} derivative DOFs, Hermitian reduced eigenproblem over $\Gamma$-X-M-$\Gamma$. \\
\textbf{One new ingredient} & Dynamics + periodicity + micro-inertia. \textbf{Nothing else added.} \\
\textbf{Headline result} & Band gaps become direction-dependent: rotating the ellipsoidal microstructure migrates the stop band from $\Gamma$--X to $\Gamma$--M. \\
\textbf{Validation} & \COR\ Five-layer validation and verification framework: Layers 1--4 external (3 published anchors 2023/2024/2026, closed-form analytics, independent method), Layer 5 = eight-test suite + mesh convergence/resolution floor. Hard gate: $\le2\%$ on Layers 1, 2a, 2b; all benchmark evidence in the manuscript. \\
\textbf{Length} & \STR\ 30--34 pp, 71 main-text + 14 appendix equations, 13 figures, 6 tables, $\approx$65 references. \\
\textbf{Safeguards v1.1} & \ADDED\ Independent reproduction protocol; validation evidence in manuscript; Appendices A--B; 8th internal test; provenance tags [C]/[A]/[S]; PCR1--PCR8 (\S10.4). \\
\textbf{Target} & Int.\ J.\ Mechanical Sciences (Q1, IF 11.4). Backups: JSV 5.8, Compos.\ Struct.\ 7.8, Appl.\ Math.\ Model.\ 5.1. \\
\textbf{Effort} & 14 weeks, 4 gates. \\
\textbf{Week 1} & \textcolor{ok}{\textbf{Already complete.}} \\
\end{xltabular}

\vspace{0.4cm}
\begin{center}
{\color{accent}\rule{0.85\textwidth}{0.8pt}}\\[0.3cm]
{\footnotesize \textbf{Blueprint LOCKED v1.2} --- 22 September 2026 \quad$\cdot$\quad
Any change to Sections 1--3 of the Lock Sheet requires re-approval before coding starts.\\[0.1cm]
Revision 1.1 added publication-readiness safeguards (\ADDED/\STR, Section~11); revision 1.2 consolidated the validation-layer numbering only (\COR, Section~12).\\[0.1cm]
Revision 1.3 (editorial only, 22 September 2026): Week-6 G2 gate wording corrected from ``all 7 internal tests pass'' to ``all 8 automated internal tests pass''. No scientific content changed.}\\[0.2cm]
{\color{accent}\rule{0.85\textwidth}{0.8pt}}
\end{center}


Revision A2 (v1.5, 24 September 2026): external-validation evidence routes --- quantitative, graphical, not validated; thresholds unchanged; see Section~13.
\end{document}
